\documentclass[a4paper,11pt]{article}% {{{
% -------------------------------------------------
% Set up the page margins.
% -------------------------------------------------
\usepackage[margin=1in]{geometry}
% \usepackage[headheight = .5in, headsep = \baselineskip, top = 1in, left = 1in, right = 1in, textwidth = 7.52 in]{geometry}
% \usepackage[text={6.5in,9.5in},centering]{geometry}
% \usepackage[margin=1in]{geometry}
% \setlength{\topmargin}{-0.25in}

% -------------------------------------------------
% Load Packages here
% -------------------------------------------------
\usepackage{amssymb, latexsym, amsmath, graphics, fullpage, epsfig, amsthm, relsize, pgf, tikz, amsfonts, makeidx, latexsym, ifthen, calc}
\usepackage[colorlinks,citecolor=blue,urlcolor=blue]{hyperref}
%\usepackage{eucal} this is a different font than \mathcal{•}
\usetikzlibrary{arrows}
\usepackage{mathrsfs}
\usepackage[shortlabels]{enumitem}
\usepackage{tikz}
\usepackage{amsmath}
\usetikzlibrary{arrows}
\usetikzlibrary{shadows.blur}
\usepackage{pgfplots}
\usepackage{mwe} % For dummy images
\usepackage{subcaption}
\usepackage{multirow}
\usepackage{dcolumn}
\newcolumntype{2}{D{.}{}{2.0}}
\usepackage{multicol}
\numberwithin{equation}{section}

% -------------------------------------------------
% check references
% -------------------------------------------------
% \usepackage{refcheck}
\usepackage[notref,notcite]{showkeys}
% \usepackage[notcite]{showkeys}
% \usepackage{showkeys}


% -------------------------------------------------
% Environments here
% -------------------------------------------------
\theoremstyle{plain}
\newtheorem{theorem}{Theorem}[section]
\newtheorem{corollary}[theorem]{Corollary}
\newtheorem{lemma}[theorem]{Lemma}
\newtheorem{proposition}[theorem]{Proposition}

\theoremstyle{definition}
\newtheorem{definition}[theorem]{Definition}
\newtheorem{hypothesis}[theorem]{Hypothesis}
\newtheorem{assumptions}[theorem]{Assumption}

\newtheorem{remark}[theorem]{Remark}
\newtheorem{conjecture}[theorem]{Conjecture}
\newtheorem{example}[theorem]{Example}
\newtheorem{notation}[theorem]{Notation}


% -------------------------------------------------
%  Define short hand symbols.
% -------------------------------------------------

\newcommand{\lip}{\textup{Lip}_b}
\newcommand{\liph}{\text{Lip}_{\hat\rho}}
\newcommand{\R}{\mathbb{R}}
\newcommand{\E}{\mathbb{E}}
\newcommand{\e}{\mathrm{e}}
\newcommand{\Norm}[1]{\left\|  #1   \right\|}
\newcommand{\rhoNorm}[1]{\left|  #1   \right|}
\newcommand{\ud}{\ensuremath{ \mathrm{d} }}

% }}}

\title{Invariant measures for the stochastic heat equation on $\R^d$ without drift term}
\author{Le Chen\\Auburn University\\\texttt{le.chen@auburn.edu}\\
	\and
		Nicholas Eisenberg\\Auburn University\\\texttt{nze0019@auburn.edu}
	}
\date{\today}
\allowdisplaybreaks

\begin{document}
\maketitle


\begin{abstract}
	This paper deals with the long term behavior of the solution to the nonlinear stochastic heat equation
	$\partial u /\partial t - \frac{1}{2}\Delta u= b(u)\dot{W}$ for a Lipschitz continuous diffusion coefficient
	$b$ which is not necessarily bounded. Using the theory of the stochastic integral laid out by John Walsh,
	we provide conditions which will guarantee the existence of an invariant measure for a broad range of initial
	conditions which includes $L^2_\rho$ as well as the Dirac delta distribution $\delta_0$.
\end{abstract}
\bigskip

\noindent{\it \noindent MSC 2010 subject classification}: 60H15, 60H07, 60F05.
\smallskip

\noindent{\it Keywords}: Stochastic heat equation, parabolic Anderson model, invariant measure, delta initial condition.
\smallskip

% \noindent{\it Running head:} Ergodicity and CLT for PAM.

\tableofcontents

\section{Introduction}
We consider the following stochastic heat equation
\begin{equation}
\begin{cases}
\dfrac{\partial u}{\partial t}(t,x)-\frac{1}{2}\Delta u(t,x) = b(x,u(t,x))\dot W(t,x) & \text{$x\in \R^d$ , $t>0$} \\
u(0,\cdot) = \mu(\cdot)  &
\end{cases}  \label{E:gspde}
\end{equation}
where $b(x,u)$ is assumed to be a globally Lipschitz continuous function in $u$ with Lipschitz constant independent of $x$.
In particular, the linear case  $b(x,u) = \lambda u$ is referred to as the \textit{parabolic Anderson model} \cite{CarMol94}.
In this equation, $\dot{W}(t,s)$ refers to a centered Gaussian noise that is white in time and homogeneously correlated in space. The initial measure $\mu \in \mathcal{M}_B$ where $\mathcal{M}_B$ is the space of deterministic and locally finite (regular) Borel measures which satisfy \eqref{A:icon} below. In particular, we investigate the conditions required to guarantee the existence of invariant measures for the solution to \eqref{E:gspde} for a more broad range of initial measures which include the Dirac delta distribution $\delta_0$ as well some weighted $L^2$ functions.

The question of the existence of invariant measures in unbounded domains with dimension $d \ge 3$ has been handled in few papers such as \cite{TessitoreZabczyk98M, AM03, MSY15, MSY16}. One thing to note is that the just mentioned papers use the theory of the stochastic integral laid out by Da Prato and Jabczyk in \cite{DZ92} while in this paper we use the theory developed by Walsh in \cite{Walsh}. The former develops a theory of integration with respect to
Hilbert-space-valued processes while the later emphasizes integration with respect to worthy martingale measures. Both
theories have their advantages and disadvantages but have been proven to be more or less equivalent in \cite{DalQuer}.

The need to prove existence of invariant measures is apparent for the fact that they are crucial in order to study the
ergodicity of dynamical systems. Da Prato and Jabczyk \cite{DZp92} established the following procedure
\begin{itemize}
	\item Show the existence of a Markovian solution to the SPDE in question in a certain function space where the
	transition group satisfies the Feller Property
	\item Show that the solution starting from a particular initial condition is bounded in probability
\end{itemize}
The first bullet point has been shown to be true \cite{DZ96} for a wider class of SPDEs which includes \eqref{E:gspde}
and in fact has also been show to be true in other various cases of interest \cite{Cer01, MS99}.
In \cite{TessitoreZabczyk98M} they follow this procedure to show the existence of an invariant measure for the solutions to
\eqref{E:gspde} in two separate cases:
\begin{enumerate}
	\item The initial measure is a function, $\varphi \in L^2_\rho \cap L^2_{\hat\rho}$ where both $\rho$ and
	$\hat\rho$ satisfy the condition that $\rho \cdot (\hat\rho)^{-1} \in L^1(\R^d)$ and the solution is
	bounded in probability.
	\item The initial measure is a constant function while assuming certain integrability properties of the spectral
	density.
\end{enumerate}
By now using the theory of Walsh as well as subsequent results proven in \cite{CH16Comparison} we are able to mimic the
proof in \cite{TessitoreZabczyk98M}, but in the language of Walsh, and show the existence of invariant measures for a
more broad group of initial conditions. However, because of the differences between the two theories of the stochastic
integral, we must trade some of previous assumptions on the spectral density for new ones which are listed below. For
example we drop hypothesis 2.1(i) as well as the assumptions from \cite[Theorem 3.3]{TessitoreZabczyk98M} for the new
assumptions listed below.

Misiats, Stanzhytski and Yip \cite{MSY15} show the existence of a invariant measure for \eqref{E:gspde} with any
$L^2(\R^d)$ initial condition but with non-zero drift term $f(x,u(t,x))$ added to the right side of the equal sign.
However their drift term $f$ and diffusion coefficient $b$ must be bounded above by a some $\varphi \in L^\infty(\R^d)
\cap L^1(\R^d)$ in the sense that
\begin{enumerate}
	\item $|f(x,0)|,\; |b(x,0)| \; \le \varphi(x)$
	\item $|f(x,u)-f(x,v)|, \; |b(x,u)-b(x,v)|\; \le L\varphi(x)|u-v|$
\end{enumerate}
In our paper we do not require such an upper bound. We also allow for any $L^2_\rho(\R^d)$ initial condition which
contains the $L^2(\R^d)$ initial conditions. However we only consider the case where the drift term $f \equiv 0$.  For
the case of unbounded coefficients, Manthey and Assing \cite{AM03} have showed the existence of an invariant measure
when considering the weight $\rho(x)=(1+|x|^2)^{-a/2},\; a>d$ and a bounded continuous initial condition. Their
diffusion coefficient follows the same growth properties as ours but they allow for non-zero drift term that follows a
"one-sided" linear growth condition. In their paper they are able to squeeze the solution to their SPDE in between two
other solutions, each satisfying a desired inequality that leads to the existence of an invariant measure. As of now,
the theory of Walsh can not be applied to such drift coefficients so it is unclear if the results in this paper may
apply to such a SPDE.

Before we state the main results of this paper we need to propose series of assumptions. First, we have the following:

\begin{assumptions}
	\label{A:General}
	\begin{enumerate}[(i)]
		\item The initial measure $\mu$ satisfies the following integrability condition:
		\begin{equation}
		\int_{\R^d}\exp\left(-a|x|^2\right)|\mu|(\ud x) < \infty \quad \text{for any $a>0$;} \label{A:icon}
		\end{equation}
		\item The spectral measure $\hat{f}$ satisfies {\it Dalang's condition}:
		\begin{equation}
		\Upsilon(\beta) := (2\pi)^{-d}\int_{\R^d}\dfrac{\hat f(\ud \xi)}{\beta + |\xi|^2} <
		\infty \quad \text{for some and therefore any $\beta>0$}  \label{A:dalcon}
		\end{equation}
		\item The diffusion coefficient $b(x,u)$ satisfies the following Lipschitz and boundedness conditions:
		\begin{equation}
		|b(x,u)-b(x,v)|< L |u-v| \quad \text{and} \quad |b(x,0)| \le L_0 \quad u,v \in \R, \; x \in \R^d \label{A:lipcon}.
		\end{equation}
	\end{enumerate}
\end{assumptions}

With Assumption \ref{A:General}, we can now state two already established results. The first result deals with the existence and uniqueness of the solution of \eqref{E:gspde}. The second result gives a bound on the moments $p^{th}$-moments of the solution, where we denote $\Norm{X}_p := (\E|X|^p)^{1/p}$ to be the $p^{th}$ moment of a random variable $X$. It should be noted the that proof of our main result "Insert main result" heavily relies on the moment bound given in Theorem \ref{T:ExUn} below.

\begin{theorem}[Theorem 1.2 \cite{CH16Comparison}]
	Under Assumption \ref{A:General}, the SPDE \eqref{E:gspde} has a unique random field solution starting from $\mu$ and the solution is $L^2(\Omega)$ continuous.
\end{theorem}

\begin{theorem}[Theorem 1.7 \cite{CH16Comparison}]
	\label{T:ExUn}
	Under Assumption \ref{A:General}, for any $t>0$ and $x \in \R^d$, the solution to \eqref{E:gspde}, $u(t,x)$,
	given by \eqref{E:gmsol} is in $L^p(\Omega)$ and for $p\ge 2$
	\begin{equation}
	\label{E:mb}
	\Norm{u(t,x)}_p \le \big[\bar{\varsigma}+\sqrt{2}(|\mu|*G(t,\cdot))(x)\big]H(t;\gamma_p)^{1/2},
	\end{equation}
	where $\bar{\varsigma}=L_0/L$, $\gamma_p = 32pL^2$ and the function $H(t;\gamma_p)$ is nondecreasing in $t$ (see \cite{CH16Comparison} for
	the concrete definition of $H$).
	Moreover, if we assume that
	\begin{equation}
	\Upsilon(0) < \infty \label{A:dalcon0}
	\end{equation}
	then this function $H(t;\gamma_p)$ is bounded in $t$. Hence, by choosing
	$C=\left(\sqrt{2}\vee \overline{\varsigma}\right)\sup_{t\ge 0} H(t;\gamma_p)^{1/2}<\infty$,
	we have that
	\begin{equation}
	\label{E:mb0}
	\Norm{u(t,x)}_p \le
	\begin{cases}
	C\left[1+\left(|\mu|*G(t,\cdot)\right)(x)\right] & \text{if $L_0> 0$,}\\
	C\left(|\mu|*G(t,\cdot)\right)(x) & \text{if $L_0= 0$}.
	\end{cases}
	\end{equation}
\end{theorem}

\subsection{Main Results}

Our first main result handles the case of the existence of an invariant measure for $L^2_{\rho}(\R^d)$ initial conditions (see Section \ref{S:weightedspace}), namely the following SPDE

\begin{equation}
\begin{cases}
\dfrac{\partial u}{\partial t}(t,x)-\frac{1}{2}\Delta_x u(t,x) = b(x,u(t,x)\dot W(t,x) & \text{$x\in \R^d$ , $t>0$}
\\ u(0,x)=\varphi(x) & \varphi(x) \in L^2_\rho(\R^d)
\end{cases}  \label{E:spde}.
\end{equation}

 Before we state the result, we need to make an additional assumption on the spectral measure.

\begin{assumptions}
			\item \begin{equation}
	\int_0^s \ud r \: (s-r)^{-2\alpha} \int_{\R^{d}} \hat f(\ud \xi) \exp\left(-(s-r)|\xi|^2 \right) < \infty \label{A:fcon}
	\end{equation}
\end{assumptions}

With this assumption now stated, we now state our first main result, Theorem \ref{T:Main}.

\begin{theorem}
	\label{T:Main}
	Assume that Assumption \ref{A:General} holds with the stronger assumption that $\Upsilon(0) <\infty$. Also assume \eqref{A:fcon}. Now for any admissible weight $\rho$, suppose there exists another admissible weight $\hat\rho$ and an element $\hat{\varphi} \in L^2_\rho(\R^d) \cap L^2_{\hat\rho}(\R^d)$ such that
	\begin{align*}
	\int_{\R^d}\dfrac{\rho(x)}{\hat\rho(x)}\ud x < \infty \quad \text{and $u^{\hat\varphi}(\cdot,x)$ is bounded in probability in $L^2_{\hat\rho}(\R^d)$.}
	\end{align*}
	Then for any $t_0 >0$ there exists an invariant measure for the laws of the solution of the SPDE \eqref{E:gspde}, $\{\mathscr L (u^{\hat\varphi}(t,\cdot))\}_{t \ge t_0}$,  in $L^2_\rho(\R^d)$ for any $L^2_{\rho}(\R^d)$ initial condition.
\end{theorem}

We now state our last assumption which will be required for the proof of one of our last main result, Theorem \eqref{T:Main2}.

\begin{assumptions} $\;$
	\begin{equation}
		\sup_{x\in \R^d}	\big(G(t,\cdot)*|\mu|)(x) \big) < K(t), \quad \forall t>0 \text{ and } K(t)\in \mathbb N \label{A:newgicon}
	\end{equation}
\end{assumptions}

\begin{remark}
	It is easy to see that Assumption \eqref{A:newgicon} is satisfied whenever $\mu=\varphi \in L^2_\rho(\R^d)$ is bounded and also when $\mu = \delta_0$ where $\delta_0$ is the Dirac delta measure.
\end{remark}

\begin{theorem}
	\label{T:Main2}
	Suppose the assumptions of Theorem \ref{T:Main} remain true and in addition assume \eqref{A:newgicon} holds true. Suppose there existed
	another admissible weight $\hat{\rho}$ such that
	\begin{align*}
	\int_{\R^d}\dfrac{\rho(x)}{\hat\rho(x)}\: \ud x < \infty.
	\end{align*}
	 Then there exists an invariant measure for the laws of the solutions to
	\eqref{E:gspde} in $L^2_\rho(\R^d)$ for $t>l>0$. In particular, for any $l>0$ there exists a measure that is
	invariant for $\{\mathscr L (u^\mu(t,\cdot))\}_{t>l}$ where $u$ is the solution to \eqref{E:gspde} with
	initial measure $\mu$.
\end{theorem}

\begin{corollary}
	\label{C:delta}
	Suppose the assumption of Theorem \ref{T:Main} remain true. Then for any $l>0$, the stochastic heat equation
	\begin{equation}
	\label{E:dspde}
	\begin{cases}
	\dfrac{\partial u}{\partial t}(t,x)-\frac{1}{2}\Delta_x u(t,x) = b(x,u(t,x)\dot W(t,x) & \text{$x\in \R^d$ , $t>0$}
	\\ u(0,\cdot) = \delta_0(\cdot)  &
	\end{cases}
	\end{equation}
	has an invariant measure in $L^2_\rho$ for the laws of their solution for $t>l$ and for any $\rho$ which
	there exists $\hat{\rho}$ such that $\rho/\hat{\rho} \in L^1(\R^d)$.
\end{corollary}

\begin{proof}
	The SPDE above satisfies \eqref{A:newgicon}.
	By Lemma \ref{L:Moments} we know that for an arbitrary $t_0>0$, the solution at $t=t_0$ from the SPDE is in $L^2_\rho \cap L^2_{\hat\rho}$. Theorem \ref{T:Main2} implies the result.
\end{proof}

\section{Preliminaries}
\subsection{Definition of the Gaussian noise and Walsh integral}

Following \cite{DalangEtc09Minicourse}, a Gaussian noise that is white in time and colored in space is an
$L^2(\Omega,\mathcal F,\mathbb P)$-valued zero mean Gaussian Processes on a complete probability space $(\Omega,\mathcal F,\mathbb P)$
\begin{align*}
\{ F(\psi) : \psi \in C_0^\infty(\R^{d+1}) \} ,
\end{align*}
such that
\begin{align*}
\E [ F(\psi)F(\varphi) ] = J(\psi,\varphi) := \int_0^\infty \ud s \int_{\R^d}
\Gamma(dx)(\psi(s,\cdot)*\tilde{\varphi}(s,\cdot))(x)
\end{align*}
where $\tilde{\varphi}(x):= \varphi(-x)$ and $\Gamma$ is a tempered measure on $\R^d$ that is nonnegative and nonnegative definite. When $\Gamma$ has density $\Gamma(\ud x) = f(x)\ud x$ then $J$ above can be written as
\begin{align*}
J(\psi,\phi) = \int_0^\infty \ud s \int_{\R^d} \ud x \int_{\R^d} \ud y \: \psi(s,x)f(x-y)\varphi(s,y).
\end{align*}
$f$ is often referred to as the \textit{spatial correlation} function of the Gaussian Noise and its Fourier transform
\begin{align*}
\hat{f}(\xi)=\mathcal Ff(\xi) = \int_{\R^d}\exp(-i\xi\cdot x)f(x) \: \ud x
\end{align*}
is often referred to as the spectral density. Now following \cite{Walsh} and \cite{DalangEtc09Minicourse} and by
denoting $\mathscr B_b(\R^d)$ as the family of bounded Borel subsets of $\R^d$, one may construct out of this a worthy
martingale measure $W=(W_t(A),t\ge 0, A \in \mathscr B_b(\R^d))$ where
\begin{align*}
W_t(A):= \lim \limits_{t\to \infty} F(\psi_n),
\end{align*}
where the limit is in $L^2(\Omega,\mathcal F,\mathbb P)$, $\psi_n \in C_0^\infty(\R^{d+1})$ and $\psi_n \downarrow
1_{[0,t]\times A}. $ With this construction, one may show that $(W_t(A),t\ge 0, A \in \mathscr B_b(\R^d))$ is a worth
martingale measure and its covariation measure $Q$ is given by
\begin{align*}
Q(A \times B \times [s,t)) = (t-s)\int_{\R^d} \ud x \int_{\R^d} \ud y \: 1_A(x)f(x-y)1_B(y)
\end{align*}
and its dominating measure is $K \equiv Q$. The main point is that $F$ and $W$ are related by
\begin{align*}
F(\psi) = \int_0^\infty \int_{\R^d} \psi (t,x) W(\ud t, \ud x)
\end{align*}
where the integral is a martingale measure stochastic integral as in Walsh with filtration given by
\begin{align*}
\mathcal F_t = \sigma \big(W_s(A), s \le t, A \in \mathscr B_b(\R^d) \big) \vee \mathscr N
\end{align*}
where $\mathscr N$ is the sigma field generated by the $\mathbb P$-null sets.

Now for any adapted, jointly measurable processes (with respect to $\mathcal B((0,\infty)\times \R^d)\times \mathcal F)$
random field $\{X(t,x): t>0 \text{ and } x \in \R^d\}$ such that the for all integers $p \ge 2$,
\begin{align*}
\int_0^\infty\ud s\iint_{\R^{2d}} \ud x \ud y \: \Vert X(s,y)X(s,x) \Vert_{\frac{p}{2}}f(x-y) < \infty,
\end{align*}
the stochastic integral
\begin{align*}
\int_0^\infty \int_{\R^d} X(s,y)\: W(\ud s, \ud y)
\end{align*}
is well defined in the sense of Walsh. Throughout this paper $\Vert \cdot \Vert_p$ will denote the standard
$L^p(\Omega)$ norm, $\left(\E|\cdot|^p\right)^{\frac{1}{p}}$.

With all this said, we define the mild solution to \eqref{E:gspde} as
\begin{equation}
u(t,x)=\int_{\R^d}G(t,x-y)\mu(\ud y) + \int_0^t\int_{\R^d}b(y,u(s,y))G(t-s,x-y)W(\ud s,\ud y) \label{E:gmsol}
\end{equation}
where the and integral above is a Walsh integral and
\[
G(t,x) = (2 \pi t)^{-d/2}\exp\left( -\dfrac{|x|^2}{2t} \right).
\]
We also denote
\[
J(t,x) := \int_{\R^d}G(t,x-y)\mu(\ud y)
\]
with the distinction that
\[
J_0(t,x) := \int_{\R^d}G(t,x-y)|\mu|(\ud y)
\]
and lastly,
\[
I(t,x):= \int_0^t\int_{\R^d}b(y,u(s,y))G(t-s,x-y)W(\ud s,\ud y).
\]
\begin{definition}
	\label{D:Sol}
	A process $u=(u(t,x),(t,x)\in (0,\infty)\times \R^d)$ is called a random field solution to \eqref{E:gspde} if:
	\begin{enumerate}
		\item $u$ satisfies almost surely and for all $(t,x) \in (0,\infty)\times \R^d$
		\[
		u(t,x) = \int_{\R^d}G(t,x-y)\:\mu( \ud y) +
		\int_0^t\int_{\R^d}G(t-s,x-y)b(y,u(s,y)) \:
		W(\ud s, \ud y);
		\]
		\item $u$ is adapted, that is, for all $(t,x) \in (0,\infty)\times \R^d$, $u(t,x)$ is $\mathcal F_t$ measurable;
		\item $u$ is jointly measurable with respect to $\mathcal B((0,\infty)\times \R^d) \times \mathcal F$;
		\item $\Norm{I(t,x)}_2 < + \infty $ for all $(t,x) \in (0,\infty)\times \R^d$;
		\item $I$ is $L^2(\Omega)$-continuous, that is, the function $(t,x) \mapsto I(t,x)$ mapping $(0,\infty)\times \R^d$ into $L^2(\Omega)$ is continuous.
	\end{enumerate}
\end{definition}

\subsection{Invariant measures}
Let $(E, \mathscr E)$ denote an arbitrary measure space and $\mathcal M_1(E)$ be the set of all probability measures on $(E, \mathscr E)$. Let $P_t(x,A)$, $t\ge0$, $x\in E$, $A \in \mathscr E$ be a Markovian Transition function as in \cite{DZ96}. Associated with each Markovian transition function is a semigroup of linear operators, $P_t$ defined on the set of bounded measurable functions on $E$, $\mathcal B_b(E)$ defined by

\begin{align*}
P_t \varphi (x) = \int_E P_t(x,\ud y)\varphi(y).
\end{align*}
Lastly we say that the above semigroup of linear operators satisfies the Feller Property if for each continuous and bounded function $\varphi$ on $E$ we have that $P_t\varphi (x)$ is also a bounded and continuous function on $E$. Now consider, for any $\mu(\cdot) \in \mathcal{M}_H$, the laws of the solution $u(t,x)$ of \eqref{E:gspde}
\begin{align*}
\mathscr L (u^\mu(t,\cdot))(A) = P\{\omega \in \Omega : u^\mu(t,\cdot) \in A \quad A \in \mathscr B(L^2_\rho(\R^d))\}
\end{align*} and define $P_t(\mu,A) := \mathscr L (u^\mu(t,\cdot))(A)$. The following two definitions will define what it means for a measure to be invariant for an SPDE.

\begin{definition}
A measure $\eta \in \mathcal M_1(E)$ is \textit{invariant} for $P_t$ if
\begin{align*}
\eta(A)=\int_E P_t(x,A)\eta(dx) \quad \text{for each $t\ge 0$ and $A \in \mathcal{E}$}.
\end{align*}
\end{definition}

\begin{definition}
	A measure $\eta$ is \textit{invariant for the SPDE \eqref{E:gspde}} if $\eta$ is invariant for $P_t(\mu,A) := \mathscr L (u^\mu(t,\cdot))(A)$.
\end{definition}

\begin{remark}
	Lemma \ref{L:phi} below shows that $L^2_{\rho}(\R^d) \subset \mathcal{M}_B$.
\end{remark}

As mentioned in the introduction, we said our first goal was to establish the Feller property. However, this has been established in \cite{DZ96} and in fact they verfy this for a larger class of SPDEs. The second step was to show the solution is bounded in probability which means
\begin{align*}
\forall \epsilon >0, \; \exists R>0: \; \forall t>0, \; P[|u^\varphi(t,\cdot)|_\rho \ge R] <  \epsilon .
\end{align*}
We verify this property and deduce the existence of an invariant measure for $L^2_\rho$ initial conditions in Theorem
\ref{T:Main} and then generalize the result in Theorem \ref{T:Main2} for measure valued initial conditions, and in particular, the Dirac delta
measure $\delta_0$.

\subsection{The weighted space $L^2_\rho(\R^d)$}\label{S:weightedspace}

Let $\rho$ be some nonnegative function on $\R^d$.
We denote by $L^2_\rho(\R^d)$ the weighted Hilbert space with the inner produce $\langle f,g \rangle_\rho  :=
\int_{\R^d}f(x)g(x)\rho(x) \: \ud x$ and set $|f|_\rho^2 :=\int_{\R^d}f(x)^2\rho(x) \: \ud x$.
Following Tessitore and Zabczyk \cite{TessitoreZabczyk98M}, we assume that $\rho$ is an \textit{admissible weight}
meaning that $\rho$ is a positive, bounded and continuous function in $L^1(\R^d)$ such that for $T>0$ there exists a
constant $C_\rho(T)$ such that
\begin{equation}
\big(G(t,\cdot)*\rho(\cdot)\big)(x) \le C_\rho(T) \rho(x) \quad \text{for all $t\in[0,T]$.} \label{E:aw}
\end{equation}
Several examples of admissible weights are
\begin{enumerate}
	\item $\rho(x)=\exp(-a|x|), \; a>0$; and
	\item $\rho(x) = (1+|x|^a)^{-1}, \; a>d$.
\end{enumerate}

For the remainder of this section, we focus ourselves on \eqref{E:gspde} with initial condition $\varphi \in L^2_\rho(\R^d)$,
namely the following stochastic heat equation.
\begin{equation}
\begin{cases}
\dfrac{\partial u}{\partial t}(t,x)-\frac{1}{2}\Delta_x u(t,x) = b(x,u(t,x)\dot W(t,x) & \text{$x\in \R^d$ , $t>0$}
\\ u(0,x)=\varphi(x) & \varphi(x) \in L^2_\rho(\R^d)
\end{cases}  \label{E:spde}
\end{equation}
which has mild solution given by
\begin{equation}
u(t,x)=\int_{\R^d}G(t,x-y)\varphi(y)\ud y + \int_0^t\int_{\R^d}b(u(s,y))G(t-s,x-y)W(\ud s,\ud y) \label{E:msol}.
\end{equation}
We will show that the solution, denoted as $u^\varphi(t,x)$, is an object in $L^2_\rho(\R^d)$ for each $t>0$. We may
denote the solution to \eqref{E:spde} as $u(t,x)$ or simply as $u$ when there is no confusion. First we state and prove
a simple lemma showing that $\varphi \in L^2_\rho$ implies condition \eqref{A:icon}.

\begin{lemma}\label{L:RhoDom} Let $\rho$ be an admissible weight. Then we have the following inequality:
	\begin{align}
	\label{E:RhoDom}
	\exp\left(-a|x|^2\right)\le K_a\: \rho(x),\qquad \text{for all $x\in\R^d$.}
	\end{align}
\end{lemma}

\begin{proof}
The admissible weight property \eqref{E:aw} tells us that
	\begin{equation} \label{E:awbd1}
		\lim_{|x| \to \infty}\frac{\left(G(t,\cdot)*\rho(\cdot)\right)(x)}{\rho(x)} \le 		C_{\rho}(T) ,\quad \forall t \in [0,T].
	\end{equation}
We note that in order to prove \eqref{E:RhoDom}, it suffices to show that there exists some $K'_a \in \mathbb{N}$ such that
	\begin{equation} \label{E:awbd2}
	\lim_{|x| \to \infty}\frac{\exp\left(-a|x|^2\right)}{\rho(x)} \le K'_a.
	\end{equation}
	Indeed, this would imply that there would exist an $R_\epsilon$ such that
	\[
	\exp\left(-a|x|^2\right) \le (K'_a+\epsilon)\rho(x) \quad |x|>R_\epsilon
	\]
	and by the positivity and continuity of both $	\exp\left(-a|x|^2\right)$ and $\rho(x)$, we know that there exists some positive constant $K'(a,\epsilon)$ such that
	\[
	\exp\left(-a|x|^2\right) \le K'(a,\epsilon)\rho(x) \quad |x|\le R_\epsilon
	\]
	and by then setting $K_a := K(a,\epsilon) = \max\{K'(a,\epsilon),(K_a'+\epsilon)\}$ then
	\[
	\exp\left(-a|x|^2\right) \le K_a\rho(x) \quad x \in \R^d
	\]
	which is precisely \eqref{E:RhoDom}.

	In order to show \eqref{E:awbd2}, it suffices to show that
	\begin{equation} \label{E:awbd3}
	\lim_{|x| \to \infty} \frac{\left(G(t,\cdot)*\rho(\cdot)\right)(x)}{\exp\left(-a|x|^2\right)} = \infty
	\end{equation}
	because this would mean that $\exp\left(-a|x|^2\right)$ is much smaller than $\left(G(t,\cdot)*\rho(\cdot)\right)(x)$ for all large $|x|$ and since \eqref{E:awbd1} is true then this would mean that \eqref{E:awbd2} would be true as well by squeeze theorem,
	\[
	0 \le \lim_{|x| \to \infty}\frac{\exp\left(-a|x|^2\right)}{\rho(x)} \le  \lim_{|x| \to \infty} \frac{\left(G(t,\cdot)*\rho(\cdot)\right)(x)}{\rho(x)} \le C_{\rho}(T)
	\]

	We prove \eqref{E:awbd3} as follows: first recall $|\rho(x)| < M$ for some $M \in \mathbb N$. Then,
	\begin{align*}
	0 \le	\dfrac{\int_{\mathbb R^d}\exp\left(-\frac{|x-y|^2}{2t}\right)\rho(y)\ud y}{\exp\left(-a|x|^2\right)} &\le \dfrac{M\int_{\mathbb R^d}\exp\left(-\frac{|x-y|^2}{2t}\right) \ud y}{\exp\left(-a|x|^2\right)}
	\\&  = \dfrac{M*\sqrt{2 \pi t}}{\exp\left(-a|x|^2\right)}
	\\& \to \infty \quad \text{as } |x|\to \infty
	\end{align*}

	Thus \eqref{E:awbd3} is proved which means that \eqref{E:awbd2} is also true.

\end{proof}

\begin{lemma}
	\label{L:phi}
	Suppose that $\varphi \in L^2_\rho(\R^d)$. Then $\varphi$ satisfies \eqref{A:icon}.
\end{lemma}

\begin{proof}
By Lemma \ref{E:RhoDom} and the Cauchy-Schwatz inequality, there exists a constant $K$ such that
	\begin{align*}
	\int_{\R^d} \e^{-a|x|^2} |\varphi(x)|\ud x
	&\le K\left(\int_{\R^d} \e^{-a|x|^2} |\varphi(x)|^2\ud x\right)^{1/2}
	\le K \Norm{\varphi}_\rho<\infty.
	\end{align*}
\end{proof}

Lemma \ref{L:phi} and Theorem \ref{T:ExUn} implies that \eqref{E:spde} has a unique solution.
The following theorem tells us that the solution for any $t>0$, $u^{\varphi}(t,\cdot)$, is a element in $L_\rho^2(\R^d)$.

\begin{theorem}
	\label{T:Moment}
	Suppose Assumption \ref{A:General} holds.
	Let $u(t,x)$ be the unique solution to \eqref{E:spde} starting from  $\varphi(\cdot) \in L^2_\rho$.
	Then for all $T>0$ and for all $t \in [0,T]$, $u(t,\cdot)\in L^2_\rho(\R^d)$ with
	\begin{align}
	\label{E:uNormRho}
	\E\left(| u(t,\cdot)|_\rho^2\right)
	\le  2 C_T \left[ \Norm{\rho}_{L^1(\R^d)} + C_\rho(T) |\varphi|_{\rho}^2 \right]<\infty.
	\end{align}
	where we use $C_T:= H(T,\gamma_2)$. Moreover, if $\Upsilon(0)<\infty$ and assume \eqref{A:newgicon}, in other words, for each $t>0$ there exists a $K(t) \in \mathbb N$ such that $\sup\limits_{x\in \R^d}\big(|\varphi|*G(t,\cdot)\big)(x) < K(t)$, then one can show
	\begin{align*}
	\sup_{t\ge 0}\E\left(|u(t,\cdot)|_\rho^2\right) < \infty
	\end{align*}
\end{theorem}

\begin{proof}
	Fix an arbitrary $T>0$ and denote $C_T:=H(T,\gamma_2)^{	1/2}$. For any $t\in [0,T]$, by Fubini's Theorem and the moment bound \eqref{E:mb}, we see that
	\begin{align*}
	\E\left(| u(t,\cdot)|_\rho^2\right)
	&= \E \left(\int_{\R^d}\vert u(t,x) \vert^2 \rho(x)\ud x\right)
	\le C_T \int_{\R^d}\bigg(1+(|\varphi|*G(t,\cdot))(x)\bigg)^2\rho(x)\ud x.
	\end{align*}
	By definition $\rho\in L^1(\R^d)$, hence we see that
	\begin{align}
	\E\left(| u(t,\cdot)|_\rho^2\right)
	\le 2 C_T \left[ \Norm{\rho}_{L^1(\R^d)} + \int_{\R^d}\left(\int_{\R^d}G(t,x-y)|\varphi(y)|dy\right)^2\rho(x)\ud x\right]. \label{mstep}
	\end{align}
	Another application of the Cauchy-Schwatz inequality shows that
	\begin{align*}
	\E\left(| u(t,\cdot)|_\rho^2\right)
	\le 2 C_T \left[ \Norm{\rho}_{L^1(\R^d)} + \int_{\R^d}(\rho*G(t,\cdot))(y)|\varphi(y)|^2\ud y \right].
	\end{align*}
	Then an application of the admissible weight property \eqref{E:aw} proves \eqref{E:uNormRho}.

	% 	Lemma \ref{L:phi} along with Assumption \ref{A:General} guarantee the existence and uniqueness of the  solution which is show in \cite{CH16Comparison}. We only prove the finiteness of the supremum. Below we use $C_i$ to represent a constant.
	% \begin{align*}
	% \E\left(| u(t,\cdot)|_\rho^2\right) &= \E \int_{\R^d}\vert u(t,x) \vert^2 \rho(x)dx
	% \\&\le C_1 \int_{\R^d}\bigg(1+(|\varphi|*G(t,\cdot))(x)\bigg)^2\rho(x)dx  \quad \text{by the moment bounds \eqref{E:mb0}}
	% \\&\le C_2 \int_{\R^d}\bigg(1+(|\varphi|*G(t,\cdot))^2(x)\bigg)\rho(x)dx
	% \\ &= C_2\left[\int_{\R^d}\rho(x)\: \ud x +\int_{\R^d}\left(\int_{\R^d}G(t,x-y)|\varphi(y)|dy\right)^2\rho(x)dx \right]
	% \\ &\le C_3\left[1+\int_{\R^d}\left(\int_{\R^d}G(t,x-y)|\varphi(y)|dy\right)^2\rho(x)dx\right]
	% \\ &\le C_4\left[1+\int_{\R^d}\int_{\R^d}G(t,x-y)|\varphi(y)|^2\rho(x)dxdy\right]
	% \\ &= C_4 \left[1+\int_{\R^d}(\rho*G(t,\cdot))(y)|\varphi(y)|^2dy\right].
	% \end{align*}
	% Now by the admissible weight property \eqref{E:aw}, for any $t_0>$0, $t \in [0,t_0]$
	% \begin{align*}
	% \\ \E\left(| u(t,\cdot)|_\rho^2\right)&\le C_4 C(t_0)\left(1+\int_{\R^d}\rho(y)|\varphi(y)|^2dy\right)
	% \\  &< K\cdot C(t_0).
	% \end{align*}

	Hence,
	\[
	\sup_{0 \le t \le t_0}\E |u(t,\cdot)|_\rho^2 < \infty.
	\]

	Now consider $t\ge t_0>0$ and further assume that $\Upsilon(0) < \infty$ and that $\sup\limits_{x\in \R^d}(|\varphi|*G(t,\cdot))(x) \le K(t)$ for any $t>0$. Now starting with \eqref{mstep}, we see that
	\begin{align*}
		\E\left(| u(t,\cdot)|_\rho^2\right)
		& \le  C\left[1+\int_{\R^d}(G(t,\cdot)*|\varphi(\cdot)|)(x)^2\rho(x)dx\right]
	\end{align*}
	However,
	\begin{align*}
		\big(G(t,\cdot)*|\varphi(\cdot)|\big)(x) = \big(G(t-t_0,\cdot)*G(t_0,\cdot)*|\varphi(\cdot)|\big)(x) \le K(t_0) \big(G(t-t_0,\cdot)*1\big)(x) = K(t_0)
	\end{align*}
	where $K(t_0)$ is the bound for $\sup\limits_{x\in \R^d}(|\varphi|*G(t,\cdot))(x)$. Now continuing, we see that
	\begin{align*}
		\E\left(| u(t,\cdot)|_\rho^2\right)& \le  C\left[1+\int_{\R^d}K(t_0)^2\rho(x)dx\right]
																		\\ & = C \left[ 1+ K(t_0)^2\Norm{\rho}_{L^1(\R^d)}  \right] < \infty.
	\end{align*}
	This proves that
	\begin{align*}
		\sup_{t \ge t_0}\E |u(t,\cdot)|_\rho^2 < \infty
	\end{align*}
	thus completing the proof of the claim.
\end{proof}

\section{Proof of Theorems \ref{T:Main} and \ref{T:Main2}}

	We start by stating a lemma that will be used to complete the proof.
The proof of the lemma is very technical and serves no real purpose in understanding its use to complete
the proof of Theorem \ref{T:Main} For that reason we postpone proving it until later.

\begin{lemma}
	\label{L:Main}
	Under the assumptions of Theorem \ref{T:Main} and for arbitrary $\epsilon \in (0,1)$, $0< t_0 \le 1$ and $R>0$ there exists a compact set $\mathscr K$ in $L^2_{ \rho}$ such that
	\begin{equation}
	\label{E:cpt}
	\mathbb{P}\{u^{\hat{\varphi}}(t,\cdot) \in \mathscr K \} \ge (1-\epsilon)\mathbb{P}\{|u^{
		\hat{\varphi}}(t-t_0,\cdot)|_{\hat \rho} < R\} \quad \text{for all $t>t_0$}.
	\end{equation}
\end{lemma}

We may now complete the proof of Theorem \ref{T:Main}, but first we must recall two definitions.

\begin{definition}[Bonded in Probability]\label{D:bdinprob}
	The process $u^{\varphi}(t,x)$ is said to be \textit{bounded in probability} in $L^2_{\rho}(\R^d)$ if
 	\begin{align}
	 	\forall \epsilon >0, \; \exists R>0: \; \forall t>0, \; \mathbb{P}\{ |u^{\varphi}(t,\cdot)|_{\rho} \ge R \} < \epsilon.
 	\end{align}
\end{definition}

\begin{definition}[Tightness of a family of measures] \label{D:tight}
	Consider the Hausdorff space $(X,T)$ and suppose that the sigma algebra $\Sigma$ contains $T$. Let $M$ be a collection of probability measures defined on $\Sigma$. The collection $M$ is said to be \textit{tight} if for any $\epsilon>0$, there exists a compact set $K_{\epsilon} \subset X$ such that for any measure $m\in M$, $m(K_{\epsilon}) > 1 - \epsilon.$
\end{definition}

 \begin{proof}[Proof of Theorem \ref{T:Main}]
 	First recall that by Theorem \ref{T:Moment}
 		\begin{equation} \label{E:bipc}
 			\sup_{t \ge 0} \E |u^{\hat{\varphi}}(t,\cdot)|_\rho^2 < \infty.
 		 \end{equation}
 	In fact, this shows that $u^{\hat{\varphi}}(t,\cdot)$ is bonded in probability in $L^2_{\rho}(\R^d)$. To see this, let $\sup_{t\ge 0} \E |u(t,\cdot)|_\rho^2 = B$ and let $\lambda >0$.
 	Then by applying Chebyshev's inequality we may choose $C>0$ large enough so that
 	\begin{align*}
 	\mathbb{P}[|u(t,\cdot)|_\rho \ge C] \le C^{-2}\E |u(t,\cdot)|_\rho^2 \le C^{-2}B < \lambda.
 	\end{align*}
 	Hence $u(t,\cdot)$ is bounded in probability in $L^2_{\rho}$.

 	Now let $\epsilon \in (0,1)$. Since $u^{\hat\varphi}(t,x)$ is bounded in probability, we can choose $R>0$ such that for any $t \ge t_0$,
 		\[
 			\mathbb{P}\{|u^{\hat \varphi}(t-t_0,\cdot)|_{\hat \rho} \ge R\} < \epsilon.
 		\]
 	By applying Lemma \ref{L:Main}, we can find a compact set $\mathscr K \in L^2_{\rho}$ such that
 	\begin{align*}
 	\mathbb{P}\{u^{\hat \varphi}(t,\cdot) \in \mathscr K \} \ge (1-\epsilon)\mathbb{P}\{|u^{\hat \varphi}(t-t_0,\cdot)|_{\hat \rho} < R\} \quad \text{for all $t>t_0$} .
 	\end{align*}
 	In return, this implies that the family of laws $\mathscr L(u^{\hat \varphi}(t,\cdot))$ for $t \ge t_0$ is tight in $L^2_{\rho}$. Indeed, if we multiply both sides of \eqref{E:cpt} by $-1$ and then add $1$ to both sides we see that
 		\begin{align*}
 			\mathbb{P}\{u^{\hat \varphi}(t,\cdot) \not\in \mathscr K \} &\le \mathbb{P}\{|u^{\hat \varphi}(t-t_0,\cdot)|_{\hat \rho} \ge R\} + \epsilon \mathbb{P}\{ |u^{\hat \varphi}(t-t_0,\cdot)|_{\hat \rho} < R \}
 			\\& \le 2 \epsilon.
 		\end{align*}
 	Therefore, the family of probability measures
 	\begin{equation}
 	\label{E:intlaw}
 	\dfrac{1}{T}\int_{t_0}^{T+t_0} \mathscr L(u^{\hat \varphi}(t,\cdot))\: \ud t, \quad T>0
 	\end{equation}
 	is tight as well; see \cite[Corollary 3.1.2]{DZ96}.
 	We may say that \eqref{E:intlaw} has a subsequence that weakly converges to some measure $\eta$ and that in fact $\eta$
 	is invariant for the laws of \eqref{E:msol}.
 \end{proof}

We now prove a lemma which will be used to prove Theorem \ref{T:Main2}.

\begin{lemma}
	\label{L:Moments}
	Consider the stochastic heat equation \eqref{E:gspde} and suppose the assumption of Theorem \ref{T:Main} remain true and \eqref{A:newgicon} holds true. Then for any $t_0>0$, $u(t_0,\cdot)\in L^2_\rho$ almost surely for any admissible weight. Also the solution to the following stochastic heat equation
	\begin{equation}
	\label{E:vspde2}
	\begin{cases}
	\dfrac{\partial v}{\partial t}(t,x)-\frac{1}{2}\Delta_x v(t,x) = b(x,v(t,x)\dot W(t,x) & \text{$x\in \R^d$ , $t>0$}
	\\ v(0,x) = u(t_0,x)  &
	\end{cases}
	\end{equation}
	satisfies the property that
	\begin{align*}
	\sup_{t\ge 0} \E |v(t,\cdot)|_\rho^2 < \infty.
	\end{align*}
\end{lemma}

\begin{proof} We first show that $u(t_0,x)\in L^2_\rho$ for any admissible weight.
	\begin{align*}
		\E |u(t_0,\cdot)|_\rho^2 &= \int_{\R^d}\ud x \: \Norm{u(t_0,x)}_2^2\rho(x)
													 \\& \le C  \int_{\R^d}\ud x \: \big(1+(|\mu|*G(t_0,\cdot)(x))\big)^2\rho(x)
	\end{align*}
	which follows from the moment estimates \eqref{E:mb0}. Now by applying \eqref{A:newgicon} we see that
	\begin{align*}
		\E |u(t_0,\cdot)|_\rho^2 & \le C \int_{\R^d}\ud x \: \big(1+K(t_0)\big)^2\rho(x)
													 \\& = C\big(1+K(t_0)\big)^2\Norm{\rho}_{L^1(\R^d)}
	\end{align*}
	which shows that $u(t_0,x)\in L^2_\rho$ almost surely for any admissible weight. Now it remains to show $\sup_{t\ge 0} \E |v(t,\cdot)|_\rho^2 < \infty$.
	\begin{align*}
		\E |v(t,\cdot)|_\rho^2 &= \int_{\R^d}\Norm{v(t,x)}_2^2\rho(x) \: \ud x
												 \\& \le C \int_{\R^d} \left( 1+\left( \Norm{u(t_0,\cdot)}_2 * G(t,\cdot)(x) \right) \right)^2\rho(x) \: \ud x
	\end{align*}
	which follows from an extension of moment estimates \eqref{E:mb0} for random initial conditions.We now continue by again applying the moment estimates.
	\begin{align*}
		\E |v(t,\cdot)|_\rho^2 & \le C \int_{\R^d} \bigg( 1+ \big(\big[1+(\mu*G(t_0,\circ))\big](\cdot) * G(t,\cdot)\big)(x)  \bigg)^2\rho(x) \: \ud x
												 \\& = C \int_{\R^d} \bigg( 1+\bigg( 1+\big[(\mu*G(t_0,\circ))(\cdot) * G(t,\cdot)\big](x) \bigg) \bigg)^2\rho(x) \: \ud x
	\end{align*}
	Now an application of \eqref{A:newgicon} gives us
	\begin{align*}
		\E |v(t,\cdot)|_\rho^2 &\le C \int_{\R^d} \bigg( 1+\bigg( 2+K(t_0) \bigg) \bigg)^2\rho(x) \: \ud x
												 \\&= C\bigg( 3 +K(t_0) \bigg)^2 \Norm{\rho}_{L^1{\R^d}}
	\end{align*}
	which is true for every $t>0$. This shows $\sup_{t\ge 0} \E |v(t,\cdot)|_\rho^2 < \infty$.
\end{proof}

We now can prove Theorem \ref{T:Main2}.

\begin{proof}[Proof of Theorem \ref{T:Main2}] Pick $t_0>0$ and $l_0>0$ such that $l=t_0+l_0$ for a given $l >0$.  By Theorem
	\ref{T:Main}, there exists an invariant measure, $\eta$, for $\{\mathscr L(v(t,\cdot))\}_{t\ge l_0}$ where $v$ is
	the solution to \eqref{E:vspde2}. However, by a general property of the heat solution, $v(t,x)=u^\mu(t+t_0,x)$ and so
	$\eta$ is invariant for $\{\mathscr L(u^\mu(t+t_0,\cdot))\}_{t\ge l_0} = \{\mathscr L(u^\mu(t,\cdot))\}_{t\ge
		l}$.
\end{proof}

\section{Proof of Lemma \ref{L:Main}}
We start off  with a proposition whose proof is given in Proposition 2.1 \cite{TessitoreZabczyk98M}

\begin{proposition}
	\label{P:weight}
	For any admissible weight $\rho$, the operators on $L^2_\rho(\R^d)$ defined by $\varphi \mapsto \big(G(t,\cdot)*\varphi(\cdot)\big)(x)$ can be extended to a $C_0-semigroup$ on $L^2_\rho$. Moreover, if $\hat\rho$ is another admissible weight such that
	\begin{align*}
	\int_{\R^d}\dfrac{\rho(x)}{\hat\rho(x)}\ud x < \infty
	\end{align*}
	then for any $t<0$, the operators defined above are compact from $L^2_{\hat\rho}(\R^d)$ to $L^2_\rho(\R^d)$.
\end{proposition}

We now define two operators. Let $q>2$ and $t_0>0$ be arbitrary and $\alpha \in (q^{-1},2^{-1})$. Define the operator $F_\alpha^{t_0}$ as follows
\begin{equation}
\label{E:fop}
\left(F_{\alpha}^{t_0} f\right)(x) := \int_0^{t_0} ds \int_{\R^d} dy \: (t_0-s)^{\alpha-1}G(t_0-s,x-y)f(s,y) .
\end{equation}
It is shown in the proof of Theorem 3.1 \cite{TessitoreZabczyk98M} that $F_\alpha$ is compact from $L^q\big((0,t_0),L^2_{\hat\rho} \big)$ to
$L^2_\rho$ where $L^q\big((0,t_0),L^2_{\hat\rho} \big)$ is the space of functions that map $s \mapsto f(s,\cdot)$ and
$f(s,\cdot) \in L^2_{\hat\rho}(\R^d)$ and in addition $\int_0^{t_0} |f(s,\cdot)|_{\hat\rho}^q \: \ud s < \infty$.

Next we define for $\varphi \in L^2_{\rho}(\R^d)$ the operators $Y^\varphi(s,y)$, or simply as $Y(s,y)$ when there is no
confusion, as follows
\begin{equation}
\label{E:yop}
Y^\varphi(s,y) =\int_0^s \int_{\R^d} (s-r)^{-\alpha}G(s-r,y-z)b(u^\varphi(r,z))W(dr,dz).
\end{equation}

\begin{lemma}
	\label{L:yineq}
	Let $q>2$, then $Y^\varphi$ operators, \eqref{E:yop}, satisfies the inequality
	\begin{equation}
	\label{E:yineq}
	\E\int_0^1 |Y^\varphi(s,\cdot)|_{\rho}^q ds \le K|\varphi|_\rho^q .
	\end{equation}
\end{lemma}

\begin{proof} First, notice by applying Minkowski's Integral Inequality that
	\begin{align*}
	E |Y(s,\cdot)|_{\rho}^q &= \int_\Omega dP \left(\int_{\R^d} dy Y(s,y)^2\rho(y) \right)^{q/2}
	\\ &\le \left(\int_{\R^d}dy \left( \int_\Omega dP \; Y(s,y)^q \rho(y)^{q/2} \right)^{2/q}\right)^{q/2}
	\\ &= \left(\int_{\R^d}dy \rho(y) \Norm{ Y(s,y) }_q^2 \right)^{q/2}
	\\ &= \big| \Norm{ Y(s,\cdot) }_q \big|_\rho^q
	\end{align*}
	and because of this, we look to find an appropriate upper bound for $\Norm{Y(s,y)}^2_q$ which for now we denote as
	$M(s,y)$. Then we integrate this bound $$ \int_{\R^d}M(s,y)\rho(y)\:\ud y$$ which gives us a bound for $\big|
	\Norm{Y(s,y)}_q\big|_\rho$. To do this, we start by applying the Burkholder-Davis-Gundy inequality, Theorem 5.27
	\cite{DalangEtc09Minicourse}, as well as Minkowski's integral inequality to see that
	\begin{align*}
	\Norm{Y(s,y)}_q^2 \le & C_q \int_0^s \ud r \: (s-r)^{-2\alpha} \iint_{\R^{2d}} \ud z_1\ud z_2 G(s-r,y-z_1)
	\Norm{b(u(r,z_1))}_q f(z_1-z_2)\\
	&\hspace{13.5em} \times  G(s-r,y-z_2) \Norm{b(u(r,z_2))}_q
	\end{align*}
	where $C_q$ is the standard constant from the BDG-inequality. Note that for Lipschitz functions, we have that
	$|b(u)|<k(1+|u|)$. We apply this and the moment bound \eqref{E:mb0} to $\Norm{b(u(r,z_i))}_q$ above to see that
	\begin{align*}
	\Norm{b(u(r,z_i))}_q
	& \le k(1+\Norm{u(r,z_i)}_q)
	\\ & \le k\bigg(1+(1+J_0(r,z_i))\bigg) \quad \text{by \eqref{E:mb0}}
	\\ & \le k(1+J_0(r,z_i)).
	\end{align*}
	Now applying this above  to the inequality for $\Norm{Y(s,y)}_q^2$ we get that

	\begin{align*}
	\Norm{Y(s,y)}_q^2 &\le k \int_0^s \ud r \: (s-r)^{-2\alpha} \iint_{\R^{2d}} \ud z_1\ud z_2 f(z_1-z_2) \prod_{i=1}^2  G(s-r,y-z_i)  (1+J_0(r,z_i))
	\\ &=  k \int_0^s \ud r \: (s-r)^{-2\alpha} \iint_{\R^{2d}} \ud z_1\ud z_2 G(s-r,y-z_1)f(z_1-z_2)G(s-r,y-z_2)
	\\ &\hspace{1.5em}+  k \int_0^s \ud r \: (s-r)^{-2\alpha} \iint_{\R^{2d}} \ud z_1\ud z_2 G(s-r,y-z_1)f(z_1-z_2)G(s-r,y-z_2)J_0(r,z_1)
	\\ &\hspace{1.5em}+  k \int_0^s \ud r \: (s-r)^{-2\alpha} \iint_{\R^{2d}} \ud z_1\ud z_2 G(s-r,y-z_1)f(z_1-z_2)G(s-r,y-z_2)J_0(r,z_2)
	\\ &\hspace{1.5em}+  k \int_0^s \ud r \: (s-r)^{-2\alpha} \iint_{\R^{2d}} \ud z_1\ud z_2 G(s-r,y-z_1)f(z_1-z_2)G(s-r,y-z_2)J_0(r,z_1)J_0(r,z_2)\\
	&=: I_1 +I_2+I_3+I_4.
	\end{align*}

	We now investigate each of the integrals $I_k$ above. Throughout the remainder of this proof and the proof of Lemma 4.2 we will be using the following relation
	\begin{equation}
	\label{E:gsim}
	G(s-r,y-z_i)G(r,z_i-\sigma)=G(s,y-\sigma)G\left(\dfrac{r(s-r)}{s},z_i-\sigma-\dfrac{r}{s}(y-\sigma)\right) .
	\end{equation}

	\noindent \textbf{For $I_1$:}
	\begin{align*}
	I_1 & = \int_0^s \ud r \: (s-r)^{-2\alpha} \iint_{\R^{2d}} \ud z_1\ud z_2 G(s-r,y-z_1)f(z_1-z_2)G(s-r,y-z_2)
	\\  & = \int_0^s \ud r \: (s-r)^{-2\alpha} \int_{\R^{d}} \Gamma(\ud \xi) \big(G(s-r,y-\cdot)*G(s-r,y+\cdot)\big)(\xi).
	\end{align*}
	Now recalling that $\Gamma$ has density $f$ and by applying Plancherel's theorem, $I_1$ is bounded above by
	\begin{align*}
	(2\pi)^{-2d}\int_0^s \ud r \: (s-r)^{-2\alpha} \int_{\R^{d}} \hat f(\ud \xi) \exp\left(-(s-r)|\xi|^2 \right).
	\end{align*}
	Now if we integrate this in $y$ over $\R^d$ with respect to $\rho(y)\ud y$ we get that
	\begin{align*}
	\int_{\R^d}I_1\cdot \rho(y)\ud y &\le (2\pi)^{-2d} \int_0^s \ud r \: (s-r)^{-2\alpha} \int_{\R^{d}} \hat f(\ud \xi) \exp\left(-(s-r)|\xi|^2 \right)  \int_{\R^d}\rho(y)\ud y
	\\ &\le (2\pi)^{-2d}\cdot \Norm{\rho}_{L^1(\R^d)} \int_0^s \ud r \: (s-r)^{-2\alpha} \int_{\R^{d}} \hat f(\ud \xi) \exp\left(-(s-r)|\xi|^2 \right) .
	\end{align*}


	\noindent \textbf{For $I_2$:}
	\begin{align*}
	I_2 &= \int_0^s \ud r \: (s-r)^{-2\alpha} \iint_{\R^{2d}} \ud z_1\ud z_2 G(s-r,y-z_1)f(z_1-z_2)G(s-r,y-z_2)J_0(r,z_1)
	\\&= \int_0^s \ud r \: (s-r)^{-2\alpha} \iint_{\R^{2d}} \ud z_1\ud z_2 G(s-r,y-z_1)f(z_1-z_2)G(s-r,y-z_2)\int_{\R^d}\ud\sigma \: G(r,z_1-\sigma)|\varphi(\sigma)|
	\\ &=  \int_0^s \ud r \: (s-r)^{-2\alpha} \iint_{\R^{2d}} \ud z_1\ud z_2 G(s-r,y-z_2)f(z_1-z_2)G\left(\dfrac{r(s-r)}{s},z_1-\sigma-\dfrac{r}{s}(y-\sigma)\right)
	\\& \hspace{10em} \times \int_{\R^d}\ud\sigma \: G(s,y-\sigma)|\varphi(\sigma)|
	\\& \le (2\pi)^{-2d} \int_0^s \ud r \: (s-r)^{-2\alpha} \int_{\R^{d}} \hat f(\ud\xi) \exp\left( -\left(\dfrac{s-r}{2}+\dfrac{r(s-r)}{2s}\right)|\xi|^2\right)\int_{\R^d}\ud\sigma \: G(s,y-\sigma)|\varphi(\sigma)|.
	\end{align*}
	Now if we integrate this in $y$ over $\R^d$ with respect to $\rho(y)\ud y$ and apply the admissible weight property we get that
	\begin{align*}
	\int_{\R^d}I_2\cdot \rho(y)\ud y &\le (2\pi)^{-2d} \int_0^s \ud r \: (s-r)^{-2\alpha} \int_{\R^{d}} \hat f(\ud\xi) \exp\left( -\left(\dfrac{s-r}{2}+\dfrac{r(s-r)}{2s}\right)|\xi|^2\right)
	\\& \hspace{15em} \times \int_{\R^d}\ud y \: \rho(y)\int_{\R^d}\ud\sigma \: G(s,y-\sigma)|\varphi(\sigma)|
	\\& \le (2\pi)^{-2d} C(S)  \int_0^s \ud r \: (s-r)^{-2\alpha} \int_{\R^{d}} \hat f(\ud\xi) \exp\left( -\left(\dfrac{s-r}{2}+\dfrac{r(s-r)}{2s}\right)|\xi|^2\right)
	\\& \hspace{15em} \times \int_{\R^d}\ud\sigma \:\rho(\sigma) |\varphi(\sigma)|
	\\&= (2\pi)^{-2d} C(S) \Norm{\varphi}_{L^1_\rho(\R^d)} \int_0^s \ud r \: (s-r)^{-2\alpha} \int_{\R^{d}} \hat f(\ud\xi) \exp\left( -\left(\dfrac{s-r}{2}+\dfrac{r(s-r)}{2s}\right)|\xi|^2\right)
	\end{align*}
	where the above is true for any $0 \le s \le S$, $S>0$ is arbitrary and $C(S)$ is the constant from the admissible weight property.

	\noindent \textbf{For $I_3$:} By completing the same steps as for $I_2$ we get that
	\begin{align*}
	I_3 \le (2\pi)^{-2d} \int_0^s \ud r \: (s-r)^{-2\alpha} \int_{\R^{d}} \hat f(\ud\xi) \exp\left( -\left(\dfrac{s-r}{2}+\dfrac{r(s-r)}{2s}\right)|\xi|^2\right)\int_{\R^d}\ud\sigma \: G(s,y-\sigma)|\varphi(\sigma)|
	\end{align*}
	and now if we integrate this in $y$ over $\R^d$ with respect to $\rho(y)\ud y$ we will get the same bound as in for $I_2$ above. Namely,
	\begin{align*}
	\int_{\R^d}I_3\cdot \rho(y)\ud y &\le (2\pi)^{-2d} C(S) \Norm{\varphi}_{L^1_\rho(\R^d)} \int_0^s \ud r \: (s-r)^{-2\alpha} \int_{\R^{d}} \hat f(\ud\xi) \exp\left( -\left(\dfrac{s-r}{2}+\dfrac{r(s-r)}{2s}\right)|\xi|^2\right).
	\end{align*}
	\textbf{For $I_4$:}
	\begin{align*}
	I_4 &= \int_0^s \ud r \: (s-r)^{-2\alpha} \iint_{\R^{2d}} \ud z_1\ud z_2 G(s-r,y-z_1)f(z_1-z_2)G(s-r,y-z_2)J_0(r,z_1)J_0(r,z_2)
	\\ &= \int_0^s \ud r \: (s-r)^{-2\alpha} \iint_{\R^{2d}} \ud z_1\ud z_2 G(s-r,y-z_1)f(z_1-z_2)G(s-r,y-z_2)
	\\&\hspace{14em}\times \iint_{\R^{2d}} \ud \sigma_1\ud \sigma_2 \: G(r,z_1-\sigma_1)G(r,z_2-\sigma_2)|\varphi(\sigma_1)||\varphi(\sigma_2)|
	\\&= \int_0^s \ud r \: (s-r)^{-2\alpha} \iint_{\R^{2d}}\ud \sigma_1\ud \sigma_2 \: G(s,y-\sigma_1)G(s,y-\sigma_2)|\varphi(\sigma_1)||\varphi(\sigma_2)|
	\\& \hspace{14em}\times \iint_{\R^{2d}} \ud z_1\ud z_2\:  f(z_1-z_2)  \prod_{i=1}^2 G\left(\frac{r(s-r)}{s}, z_i-\sigma_i-\frac{r}{s}(y-\sigma_i) \right)
	\\ &\le (2\pi)^{-2d} \int_0^s \ud r \: (s-r)^{-2\alpha} \iint_{\R^{2d}}\ud \sigma_1\ud \sigma_2 \: G(s,y-\sigma_1)G(s,y-\sigma_2)|\varphi(\sigma_1)||\varphi(\sigma_2)|
	\\& \hspace{14em}\times \int_{\R^d} \hat{f}(\ud \xi) \exp\left(-\frac{r(s-r)}{s}|\xi|^2\right)
	\\&= (2\pi)^{-2d} \int_0^s \ud r \: (s-r)^{-2\alpha}\int_{\R^d} \hat{f}(\ud \xi)\exp\left(-\frac{r(s-r)}{s}|\xi|^2\right) \left(\int_{\R^{d}}\ud \sigma\:G(s,y-\sigma)|\varphi(\sigma)|\right)^2
	\\&\le (2\pi)^{-2d} \int_0^s \ud r \: (s-r)^{-2\alpha}\int_{\R^d} \hat{f}(\ud \xi)\exp\left(-\frac{r(s-r)}{s}|\xi|^2\right) \int_{\R^{d}}\ud \sigma\:G(s,y-\sigma)|\varphi(\sigma)|^2.
	\end{align*}
	Now if we integrate this in $y$ over $\R^d$ with respect to $\rho(y)\ud y$ we get that
	\begin{align*}
	\int_{\R^d}I_4\cdot \rho(y)\ud y &\le (2\pi)^{-2d} \int_0^s \ud r \: (s-r)^{-2\alpha}\int_{\R^d} \hat{f}(\ud \xi)\exp\left(-\frac{r(s-r)}{s}|\xi|^2\right)
	\\& \hspace{15em} \times \int_{\R^d}\ud y \: \rho(y)\int_{\R^{d}}\ud \sigma\:G(s,y-\sigma)|\varphi(\sigma)|^2
	\\ &\le (2\pi)^{-2d} C(S)\int_0^s \ud r \: (s-r)^{-2\alpha}\int_{\R^d} \hat{f}(\ud \xi)\exp\left(-\frac{r(s-r)}{s}|\xi|^2\right)\int_{\R^d}\ud \sigma \: \rho(\sigma)|\varphi(\sigma)|^2
	\\ &= (2\pi)^{-2d}  C(S)  |\varphi|_\rho^2\int_0^s \ud r \: (s-r)^{-2\alpha}\int_{\R^d} \hat{f}(\ud \xi)\exp\left(-\frac{r(s-r)}{s}|\xi|^2\right).
	\end{align*}

	\noindent Now, putting all of this together, we see that
	\begin{align*}
	\left(E |Y(s,y)|_{\rho}^q\right)^{2/q} &\le \int_{\R^d}dy \rho(y) \Norm{Y(s,y)}_q^2
	\\ &\le \int_{\R^d}dy \rho(y) (I_1+I_2+I_3+I_4)
	\\ &\le (2\pi)^{-2d} \Norm{\rho}_{L^1(\R^d)}  \int_0^s \ud r \: (s-r)^{-2\alpha} \int_{\R^{d}} \hat f(\ud \xi) \exp\left(-(s-r)|\xi|^2 \right)
	\\&\hspace{1em} + 2 (2\pi)^{-2d} C(S) \Norm{\varphi}_{L^1_\rho(\R^d)} \int_0^s \ud r \: (s-r)^{-2\alpha} \int_{\R^{d}} \hat f(\ud\xi) \exp\left( -\left(\dfrac{s-r}{2}+\dfrac{r(s-r)}{2s}\right)|\xi|^2\right)
	\\&\hspace{1em}+(2\pi)^{-2d} C(S)  |\varphi|_\rho^2\int_0^s \ud r \: (s-r)^{-2\alpha}\int_{\R^d} \hat{f}(\ud \xi)\exp\left(-\frac{r(s-r)}{s}|\xi|^2\right)
	\end{align*}
	and by assumption \ref{A:fcon} and Lemma 4.2 below,  the three integrals above are all finite any $s \ge 0$ so we may take the maximum for $0 \le s \le S$ and say there exists some constant $K_S$ dependent on $S$ such that
	\[
	E |Y(s,\cdot)|_{\rho}^q \le K_S  |\varphi|_\rho^q, \quad
	\forall s\in[0,S] .
	\]
	Lastly if we integrate out $s$ from $0\to1$ we get our desired inequality,
	\[
	\int_0^1 E |Y(s,\cdot)|_{\rho}^q ds \le \int_0^1 K_1  |\varphi|_\rho^q ds = K_1 |\varphi|_\rho^q .
	\]
\end{proof}

\begin{lemma}
	\label{L:equiv}
	Assuming \eqref{A:fcon} then we may say the following two integrals are finite
	\begin{equation}
	\int_0^s \ud r \: (s-r)^{-2\alpha} \int_{\R^d} \hat{f}(\ud \xi) \exp\left( -\dfrac{r(s-r)}{s}|\xi|^2 \right) \label{E:fcon2}
	\end{equation}
	and
	\begin{equation}
	\int_0^s \ud r \: (s-r)^{-2\alpha} \int_{\R^d} \hat{f}(\ud \xi) \exp\left[ -\left(\dfrac{s-r}{2}+\dfrac{r(s-r)}{2s}\right)|\xi|^2 \right] \label{E:fcon3}.
	\end{equation}
\end{lemma}

\begin{proof} We first focus on \eqref{E:fcon2}. By applying the substitution $u=s-r$ the we may rewrite
	\[
	\eqref{A:fcon} = \int_0^s \ud u \: u^{-2\alpha} \int_{\R^d} \hat{f}(\ud \xi) \exp\left( -u|\xi|^2 \right)
	\]
	and
	\[
	\eqref{E:fcon2} = \int_0^s \ud u \: u^{-2\alpha} \int_{\R^d} \hat{f}(\ud \xi) \exp\left( -\dfrac{u(s-u)}{s}|\xi|^2 \right).
	\]
	Since $u(s-u) \ge \dfrac{us}{2}$ for all $u\in[0,s]$ then we may say that
	\begin{align*} \eqref{E:fcon2} &\le \int_0^s \ud u \: u^{-2\alpha} \int_{\R^d} \hat{f}(\ud \xi) \exp\left( -\dfrac{u}{2}|\xi|^2 \right)
	\\ &= 2^{1-2\alpha}\int_0^{s/2} \ud t \: t^{-2\alpha} \int_{\R^d} \hat{f}(\ud \xi) \exp\left( -t|\xi|^2 \right) \quad \text{substitution $t=\dfrac{u}{2}$}
	\\& \le \eqref{A:fcon} < \infty  \end{align*}
	which completes the proof showing assumption \eqref{A:fcon} implies that \eqref{E:fcon2} is finite. The finiteness of \eqref{E:fcon3} follows easily by noting that
	\begin{align*}
	\eqref{E:fcon3} &\le \int_0^s \ud r \: (s-r)^{-2\alpha} \int_{\R^d} \hat{f}(\ud \xi) \exp\left(-\dfrac{s-r}{2}|\xi|^2 \right)
	\\ &= 2^{1-2\alpha}\int_0^{s/2} \ud u \: u^{-2\alpha} \int_{\R^d} \hat{f}(\ud \xi) \exp\left(-u|\xi|^2 \right) \quad \text{using the substitution $u=\dfrac{s-r}{2}$}
	\\ & \le \eqref{A:fcon} < \infty
	\end{align*}
	which completes the proof showing that assumption \eqref{A:fcon} implies \eqref{E:fcon3}.
\end{proof}

\begin{lemma} \textit{The following factorization holds for $\alpha \in (0,1)$
		\begin{equation}
		\dfrac{\sin(\alpha \pi)}{\pi}\int_0^t(t-s)^{\alpha-1}\left[G(t-s,\cdot)*Y(s,\cdot)\right](x) \ud s = \int_0^t  \int_{\R^d} G(t-r,x-z)b(u(r,z)) W(\ud r,\ud z) \label{E:fac}.
		\end{equation} }
\end{lemma}
\begin{proof}
	\begin{align*}
	&\dfrac{\sin(\alpha \pi)}{\pi}\int_0^t(t-s)^{\alpha-1}\left[G(t-s,\cdot)*Y(s,\cdot)\right](x) \ud s \\
	\stackrel{1}{=}&\dfrac{\sin(\alpha \pi)}{\pi}\int_0^t \ud s (t-s)^{\alpha-1}\int_{\R^d} \ud y G(t-s,x-y)
	\int_0^s \int_{\R^d} \: (s-r)^{-\alpha}G(s-r,y-z)b(u(r,z))W(\ud r,\ud z)\\
	\stackrel{2}{=}& \dfrac{\sin(\alpha \pi)}{\pi}\int_0^t \ud s (t-s)^{\alpha-1}
	\int_0^s \int_{\R^d} \: (s-r)^{-\alpha}G(t-r,x-z)b(u(r,z))W(\ud r,\ud z)\\
	\stackrel{3}{=}& \dfrac{\sin(\alpha \pi)}{\pi}
	\int_0^t  \int_{\R^d} W(\ud r,\ud z)G(t-r,x-z)b(u(r,z))
	\int_r^t \ud s\: (s-r)^{-\alpha} (t-s)^{\alpha-1} \\
	\stackrel{4}{=}& \int_0^t  \int_{\R^d} G(t-r,x-z)b(u(r,z)) W(\ud r,\ud z)
	\end{align*}
	where the last step is true provided $\alpha\in (0,1)$.
\end{proof}

\begin{remark}
	We also need to verify the use of Fubini's Theorem which is stated in stated in \cite{Walsh}. We used Fubini's Theorem twice above: first from equality 1 to 2 and then from equality 3 to 4. For reference, we will state Fubini's Theorem as given in \cite{Walsh} and verify the two uses of it in the proof of Lemma 4.2.
\end{remark}

\begin{theorem}[Fubini's Theorem] Let $(G,\mathscr G, \mu)$ be a finite measure space and $(E, \mathscr E, \nu)$ be a $\sigma-$finite measure space. Let $g(r,z,\omega,y)$, $z \in E$, $r\ge0$, $\omega \in \Omega$, $y\in G$ be a $\mathscr E \times \mathscr G$ measurable function. Suppose that
	\begin{equation} \E \left[ \int_{E\times E \times [0,T]\times G} |g(r,z_1,\omega,y)g(r, z_2,\omega,y)|K(\ud z_1. \ud z_2, \ud r)\mu(\ud y)]\right] < \infty \label{E:fcond}
	\end{equation}
	where $K$ is the dominating measure as in \cite{DalangEtc09Minicourse}. Then
	\begin{align*}
	\int_G \left[ \int_E \int_0^s g(r,z,y) W(\ud r, \ud z) \right]\mu(\ud y) = \int_E \int_0^s \left[ \int_G g(r,z,y) \mu(\ud y) \right] W(\ud r, \ud z).
	\end{align*}

\end{theorem}

\begin{remark}
	Following \cite{DalangEtc09Minicourse}, the dominating measure for this problem is given by $$K(\ud z_1. \ud z_2, \ud r)= f(z_1-z_2)\ud z_1 \ud z_2 \ud r $$ where $f$ is the correlation function of the spatially homogeneous Gaussian noise.
\end{remark}

We now verify \eqref{E:fcond} for both applications of Fubini's Theorem. \\

\noindent \textbf{From Equality 1 to 2:}

\begin{proof} We let
	\begin{align*}
	&g(z,r,\omega,y)=(s-r)^{-\alpha}G(s-r,y-z)b(u(r,z))
	\\& (G,\mathscr G, \mu)=(\R^d,\mathscr B(\R^d),G(t-s,x-y) \ud y)
	\\& (E, \mathscr E, \nu) = (\R^d, \mathscr B(\R^d), \lambda)
	\end{align*}
	and so \eqref{E:fcond} becomes
	\begin{align*}
	\eqref{E:fcond} &= \E  \int_{\R^d}\ud y \;G(t-s,x-y)\int_0^s\ud r\; (s-r)^{-2\alpha}
	\\&\hspace{8em} \times \iint_{\R^{2d}}\ud z_1 \ud z_2 \:G(s-r,y-z_1)|b(u(r,z_1))|f(z_1-z_2)G(s-r,y-z_2)|b(u(r,z_2))| .
	\end{align*}
	For simplicity, we first investigate the $\E\iint_{\R^{2d}} \ud z_1 \ud z_2$ part of the integral. We continue by applying the Lipschitz property of $b$.
	\begin{align*}
	& \E\iint_{\R^{2d}}\ud z_1 \ud z_2 \:G(s-r,y-z_1)|b(u(r,z_1)|f(z_1-z_2)G(s-r,y-z_2)|b(u(r,z_2))|
	\\ \le   & k_1   \E\iint_{\R^{2d}}\ud z_1 \ud z_2 \:G(s-r,y-z_1)(1+u(r,z_1))f(z_1-z_2)G(s-r,y-z_2)(1+u(r,z_2))
	\\ \le   & k_1  \iint_{\R^{2d}}\ud z_1 \ud z_2 \:G(s-r,y-z_1)f(z_1-z_2)G(s-r,y-z_2)
	\\       & +k_1 \iint_{\R^{2d}}\ud z_1 \ud z_2 \:G(s-r,y-z_1)\Norm{u(r,z_1)}_1f(z_1-z_2)G(s-r,y-z_2)
	\\       & +k_1 \iint_{\R^{2d}}\ud z_1 \ud z_2 \:G(s-r,y-z_1)f(z_1-z_2)G(s-r,y-z_2)\Norm{u(r,z_2)}_1
	\\       & +k_1 \iint_{\R^{2d}}\ud z_1 \ud z_2 \:G(s-r,y-z_1)\Norm{u(r,z_1)}_2f(z_1-z_2)G(s-r,y-z_2)\Norm{u(r,z_2)}_2
	\\    =: & L_1 +\cdots +L_4.
	\end{align*}
	\textbf{For $L_1$:} If we back at $I_1$ from the proof of Lemma \ref{L:yineq} then we see that
	\begin{align*}
	L_1 &\le (2\pi)^{-2d}\int_{\R^d}\hat f (\ud \xi)\: \exp(-(s-r)|\xi|^2)
	\end{align*}
	and now if we include the $\int_{\R^d}\ud y \;G(t-s,x-y)\int_0^s\ud r\; (s-r)^{-2\alpha}$ part of the integral we get that
	\begin{align*}
	\int_{\R^d}\ud y \;G(t-s,x-y)\int_0^s\ud r\; (s-r)^{-2\alpha} L_1 \le (2\pi)^{-2d} \int_0^s\ud r\; (s-r)^{-2\alpha}\int_{\R^d}\hat f (\ud \xi)\: \exp(-(s-r)|\xi|^2).
	\end{align*}

	\noindent \textbf{For $L_2$:} We first use the moment bounds from \cite{CH16Comparison} along with the fact that $\Norm{u(r,z_1)}_1 \le \Norm{u(r,z_1)}_2$.
	\begin{align*}
	L_2 & \le k\iint_{\R^{2d}}\ud z_1 \ud z_2 \:G(s-r,y-z_1)(1+J_0(r,z_1))f(z_1-z_2)G(s-r,y-z_2)
	\\ &= kL_1 + \iint_{\R^{2d}}\ud z_1 \ud z_2 \:G(s-r,y-z_1)J_0(r,z_1)f(z_1-z_2)G(s-r,y-z_2).
	\end{align*}
	But this last integral above is the same thing as $I_2$ from Corollary 4.1 so we see that
	\begin{align*}
	L_2 &\le kL_1 + (2\pi)^{-2d}\int_{\R^d}\hat f(\ud \xi) \: \exp\left(-\left(\dfrac{s-r}{2}+\dfrac{r(s-r)}{2s}\right)\right) \int_{\R^d}\ud \sigma \: G(s,y-\sigma)|\varphi(\sigma)|
	\end{align*}
	and now if we include the $\int_{\R^d}\ud y \;G(t-s,x-y)\int_0^s\ud r\; (s-r)^{-2\alpha}$ part of the integral we get that
	\begin{align*}
	& \int_{\R^d}\ud y \;G(t-s,x-y)\int_0^s\ud r\; (s-r)^{-2\alpha} L_2
	\\ & \le k\int_{\R^d}\ud y \;G(t-s,x-y)\int_0^s\ud r\; (s-r)^{-2\alpha} L_1
	\\ & \hspace{3em}+ (2\pi)^{-2d}\int_{\R^d}\ud y \;G(t-s,x-y)\int_0^s\ud r\; (s-r)^{-2\alpha}\int_{\R^d}\hat f(\ud \xi) \: \exp\left(-\left(\dfrac{s-r}{2}+\dfrac{r(s-r)}{2s}\right)\right)
	\\ & \hspace{24.5em} \times  \int_{\R^d}\ud \sigma \: G(s,y-\sigma)|\varphi(\sigma)|
	\\ & \le k (2\pi)^{-2d} \int_0^s\ud r\; (s-r)^{-2\alpha}\int_{\R^d}\hat f (\ud \xi)\: \exp(-(s-r)|\xi|^2)
	\\ & \hspace{3em}+ (2\pi)^{-2d} \int_0^s\ud r\; (s-r)^{-2\alpha}\int_{\R^d}\hat f(\ud \xi) \: \exp\left(-\left(\dfrac{s-r}{2}+\dfrac{r(s-r)}{2s}\right)\right)
	\\ & \hspace{18 em} \times\int_{\R^d}\ud y \;G(t-s,x-y) \int_{\R^d}\ud \sigma \: G(s,y-\sigma)|\varphi(\sigma)|
	\\ & = k  (2\pi)^{-2d} \int_0^s\ud r\; (s-r)^{-2\alpha}\int_{\R^d}\hat f (\ud \xi)\: \exp(-(s-r)|\xi|^2)
	\\ & \hspace{3em} +(2\pi)^{-2d}\int_0^s\ud r\; (s-r)^{-2\alpha}\int_{\R^d}\hat f(\ud \xi) \: \exp\left(-\left(\dfrac{s-r}{2}+\dfrac{r(s-r)}{2s}\right)\right)\int_{\R^d}\ud \sigma \: G(t,x-\sigma)|\varphi(\sigma)|
	\\ & \le k  (2\pi)^{-2d}\int_0^s\ud r\; (s-r)^{-2\alpha}\int_{\R^d}\hat f (\ud \xi)\: \exp(-(s-r)|\xi|^2)
	\\ & \hspace{3em}+ (2\pi)^{-2d}  \Norm{G(t,x-\cdot)\varphi(\cdot)}_{L^1(\R^d)}\int_0^s\ud r\; (s-r)^{-2\alpha}\int_{\R^d}\hat f(\ud \xi) \: \exp\left(-\left(\dfrac{s-r}{2}+\dfrac{r(s-r)}{2s}\right)\right).
	\end{align*}
	\noindent \textbf{For $L_3$:} Following the same steps for $L_2$ will get you the same exact bound as for $L_3$. \\

	\noindent \textbf{For $L_4$:} We continue as we did for $L_2$ to see that
	\begin{align*}
	L_4 &\le k\iint_{\R^{2d}}\ud z_1 \ud z_2\:G(s-r,y-z_1)(1+J_0(r,z_1))f(z_1-z_2)G(s-r,y-z_2)(1+J_0(r,z_2))
	\\&= k\iint_{\R^{2d}}\ud z_1 \ud z_2\:G(s-r,y-z_1)f(z_1-z_2)G(s-r,y-z_2)
	\\&\hspace{5em}+k\iint_{\R^{2d}}\ud z_1 \ud z_2\:G(s-r,y-z_1)J_0(r,z_1)f(z_1-z_2)G(s-r,y-z_2)
	\\&\hspace{5em}+ k\iint_{\R^{2d}}\ud z_1 \ud z_2\:G(s-r,y-z_1)f(z_1-z_2)G(s-r,y-z_2)J_0(r,z_2)
	\\&\hspace{5em}+k\iint_{\R^{2d}}\ud z_1 \ud z_2\:G(s-r,y-z_1)J_0(r,z_1)f(z_1-z_2)G(s-r,y-z_2)J_0(r,z_2).
	\end{align*}
	These first three integrals have already been handled above in $L_1, L_2$ and $L_3$ of this proof so we only need to proceed with the fourth integral. However, notice from $I_4$ above in the proof of Lemma \ref{L:yineq} that this fourth integral is bounded by the following
	\begin{align*}
	&\iint_{\R^{2d}}\ud z_1 \ud z_2\:G(s-r,y-z_1)J_0(r,z_1)f(z_1-z_2)G(s-r,y-z_2)J_0(r,z_2)
	\\ &\le (2\pi)^{-2d}\int_{\R^d} \hat{f}(\ud \xi)\exp\left(-\frac{r(s-r)}{s}|\xi|^2\right) \int_{\R^{d}}\ud \sigma\:G(s,y-\sigma)|\varphi(\sigma)|^2
	\end{align*}
	and now if we include the $\int_{\R^d}\ud y \;G(t-s,x-y)\int_0^s\ud r\; (s-r)^{-2\alpha}$ part of the integral we get that
	\begin{align*}
	&\int_{\R^d}\ud y \;G(t-s,x-y)\int_0^s\ud r\; (s-r)^{-2\alpha}\int_{\R^d} \hat{f}(\ud \xi)\exp\left(-\frac{r(s-r)}{s}|\xi|^2\right) \int_{\R^{d}}\ud \sigma\:G(s,y-\sigma)|\varphi(\sigma)|^2
	\\&= \int_0^s\ud r\; (s-r)^{-2\alpha}\int_{\R^d} \hat{f}(\ud \xi)\exp\left(-\frac{r(s-r)}{s}|\xi|^2\right)\int_{\R^d}\ud y \;G(t-s,x-y) \int_{\R^{d}}\ud \sigma\:G(s,y-\sigma)|\varphi(\sigma)|^2
	\\&= \int_0^s\ud r\; (s-r)^{-2\alpha}\int_{\R^d} \hat{f}(\ud \xi)\exp\left(-\frac{r(s-r)}{s}|\xi|^2\right)\int_{\R^{d}}\ud \sigma\:G(t,x-\sigma)|\varphi(\sigma)|^2
	\\&\le \Norm{G(t,x-\cdot)\varphi(\cdot)^2}_{L^1(\R^d)} \int_0^s\ud r\; (s-r)^{-2\alpha}\int_{\R^d} \hat{f}(\ud \xi)\exp\left(-\frac{r(s-r)}{s}|\xi|^2\right)
	\end{align*}
	and now putting this all together for $L_4$ we see that
	\begin{align*}
	\int_{\R^d}& \ud y \;G_{t-s}(x-y)\int_0^s\ud r\; (s-r)^{-2\alpha}L_4
	\\ &\le k (2\pi)^{-2d}\int_0^s\ud r\; (s-r)^{-2\alpha}\int_{\R^d}\hat f (\ud \xi)\: \exp(-(s-r)|\xi|^2)
	\\&\hspace{2em}+ 2 k  (2\pi)^{-2d} \Norm{G(t,x-\cdot)\varphi(\cdot)}_{L^1(\R^d)}\int_0^s\ud r\; (s-r)^{-2\alpha}\int_{\R^d}\hat f(\ud \xi) \: \exp\left(-\left(\dfrac{s-r}{2}+\dfrac{r(s-r)}{2s}\right)\right)
	\\&\hspace{2em} + k (2\pi)^{-2d}  \Norm{G(t,x-\cdot)\varphi(\cdot)^2}_{L^1(\R^d)} \int_0^s\ud r\; (s-r)^{-2\alpha}\int_{\R^d} \hat{f}(\ud \xi)\exp\left(-\frac{r(s-r)}{s}|\xi|^2\right).
	\end{align*}
	Considering the above arguments we may say that theres is a constant $K_{x,t}$ depending on $x$ and $t$ such that
	\begin{align*}
	\eqref{E:fcond} &\le K_{x,t}\int_0^s\ud r\; (s-r)^{-2\alpha}\int_{\R^d}\hat f (\ud \xi)\: \exp(-(s-r)|\xi|^2)
	\\&\hspace{3em}+ K_{x,t}\int_0^s\ud r\; (s-r)^{-2\alpha}\int_{\R^d}\hat f(\ud \xi) \: \exp\left(-\left(\dfrac{s-r}{2}+\dfrac{r(s-r)}{2s}\right)\right)
	\\&\hspace{3em} + K_{x,t}\int_0^s\ud r\; (s-r)^{-2\alpha}\int_{\R^d} \hat{f}(\ud \xi)\exp\left(-\frac{r(s-r)}{s}|\xi|^2\right) < \infty
	\end{align*}
	where the finiteness follows by \eqref{A:fcon} and Lemma \ref{L:equiv}. This completes the justification for the first use of Fubini's Theorem. \\
\end{proof}

\noindent \textbf{From Equality 2 to 3:}
\begin{proof} We let
	\begin{align*}
	&g(r,z,\omega,y)=(s-r)^{-\alpha}G(t-r,x-z)b(u(r,z))
	\\& (G,\mathscr G, \mu)=((0,t),\mathscr B(0,t),(t-s)^{\alpha-1} \ud s)
	\\& (E, \mathscr E, \nu) = (\R^d, \mathscr B(\R^d), \lambda)
	\end{align*}
	and so \eqref{E:fcond} becomes
	\begin{align*}
	\eqref{E:fcond} = \E \int_0^t \ud s \: (t-s)^{\alpha-1}\int_0^s \ud r \: (s-r)^{-2\alpha} \iint_{\R^{2d}} \ud z_1 \ud z_2 \: f(z_1-z_2) \prod_{i=1}^{2} G(t-r,x-z_i)|b(u(r,z_i))|
	\end{align*}
	and again for simplicity we will first investigate the $\E\iint_{\R^{2d}}\ud z_1 \ud z_2$ integral. Using the Lipschitz property, we similarly get that
	\begin{align*}
	& \E\iint_{\R^{2d}}\ud z_1 \ud z_2 \:G(t-r,x-z_1)|b(u(r,z_1))|f(z_1-z_2)G(t-r,x-z_2)|b(u(r,z_2))|
	\\ \le & k  \E\iint_{\R^{2d}}\ud z_1 \ud z_2 \:G(t-r,x-z_1)(1+u(r,z_1))f(z_1-z_2)G(t-r,x-z_2)(1+u(r,z_2))
	\\ \le & k \iint_{\R^{2d}}\ud z_1 \ud z_2 \:G(t-r,x-z_1)f(z_1-z_2)G(t-r,x-z_2)
	\\     &+k \iint_{\R^{2d}}\ud z_1 \ud z_2 \:G(t-r,x-z_1)\Norm{u(r,z_1)}_1f(z_1-z_2)G(t-r,x-z_2)
	\\     &+k  \iint_{\R^{2d}}\ud z_1 \ud z_2 \:G(t-r,x-z_1)f(z_1-z_2)G(t-r,x-z_2)\Norm{u(r,z_2)}_1
	\\     &+k  \iint_{\R^{2d}}\ud z_1 \ud z_2 \:G(t-r,x-z_1)\Norm{u(r,z_1)}_2f(z_1-z_2)G(t-r,x-z_2)\Norm{u(r,z_2)}_2
	\\=:   & K_1+\cdots+K_4.
	\end{align*}
	We apply similar techniques as in for the first justification of Fubini's Theorem to get that

	\noindent \textbf{For $K_1$:}
	\begin{align*}
	\int_0^t \ud s \: (t-s)^{\alpha-1}\int_0^s \ud r \: (s-r)^{-2\alpha} K_1 &\le \int_0^t \ud s \: (t-s)^{\alpha-1}\int_0^s \ud r \: (s-r)^{-2\alpha}
	\int_{\R^d}\hat f(\ud \xi)\: \exp\left(-(t-r)|\xi|^2\right)
	\\& \le k_{t,1}\int_0^t \ud s \: (t-s)^{\alpha-1}
	\end{align*}
	where
	\[
	k_{t,1} = \max_{0\le s \le t} \int_0^s \ud r \: (s-r)^{-2\alpha}
	\int_{\R^d}\hat f(\ud \xi)\: \exp\left(-(t-r)|\xi|^2\right)
	\]
	and in fact $k_{t,1}$ is finite for each $t>0$ by Lemma \ref{L:equiv}.

	\noindent \textbf{For $K_2$ and $K_3$:}
	\begin{align*}
	& \int_0^t \ud s \: (t-s)^{\alpha-1}\int_0^s \ud r \: (s-r)^{-2\alpha} K_2
	\\ & \le k\int_0^t \ud s \: (t-s)^{\alpha-1}\int_0^s \ud r \: (s-r)^{-2\alpha}\int_{\R^d}\hat f(\ud \xi)\: \exp\left(-(t-r)|\xi|^2\right)
	\\ & \hspace{2em}+ k\int_0^t \ud s \: (t-s)^{\alpha-1}\int_0^s \ud r \: (s-r)^{-2\alpha}\int_{\R^d}\hat f(\ud \xi)\: \exp\left(-\left(\dfrac{t-r}{2}+\dfrac{r(t-r)}{2t}\right)|\xi|^2\right)
	\\ & \hspace{20em} \times \int_{\R^d}\ud \sigma \:G(t,x-\sigma)|\varphi(\sigma)|
	\\ & \le \big(k_{t,1}  + k J_0(t,x) k_{t,2} \big) \int_0^t \ud s \: (t-s)^{\alpha-1}
	\end{align*}
	where
	\[
	k_{t,2} =  \max_{0\le s \le t} \int_0^s \ud r \: (s-r)^{-2\alpha}\int_{\R^d}\hat f(\ud \xi)\: \exp\left(-\left(\dfrac{t-r}{2}+\dfrac{r(t-r)}{2t}\right)|\xi|^2\right)
	\]
	which is finite by Lemma \ref{L:equiv}.

	\noindent \textbf{For $K_4$:}

	\begin{align*}
	&\int_0^t \ud s \: (t-s)^{\alpha-1}\int_0^s \ud r \: (s-r)^{-2\alpha} K_4
	\\&\le k\int_0^t \ud s \: (t-s)^{\alpha-1}\int_0^s \ud r \: (s-r)^{-2\alpha}\int_{\R^d}\hat f(\ud \xi)\: \exp\left(-(t-r)|\xi|^2\right)
	| \\& \hspace{2em} + 2k\int_0^t \ud s \: (t-s)^{\alpha-1}\int_0^s \ud r \: (s-r)^{-2\alpha}\int_{\R^d}\hat f(\ud \xi)\: \exp\left(-\left(\dfrac{t-r}{2}+\dfrac{r(t-r)}{2t}\right) | \xi | ^2\right)
	\\& \hspace{22em} \times \int_{\R^d}\ud \sigma \:G(t , x-\sigma) | \varphi(\sigma) |
	\\& \hspace{2em}+  k\int_0^t \ud s \: (t-s)^{\alpha-1}\int_0^s \ud r \: (s-r)^{-2\alpha}\int_{\R^d}\hat f(\ud \xi)\: \exp\left(-\dfrac{r(t-r)}{t}|\xi|^2\right)\left(\int_{\R^{d}}\ud \sigma \:G(t,x-\sigma)|\varphi(\sigma)\right)^2
	\\&\le k\big(k_{t,1}+ 2k_{t,2}J_0(t,x)+k_{t,3}J_0(t,x)^2 \big)\int_0^t \ud s \: (t-s)^{\alpha-1}
	\end{align*}
	where
	\[
	k_{t,3} = \max_{0\le s \le t} \int_{\R^d}\hat f(\ud \xi)\: \exp\left(-\dfrac{r(t-r)}{t}|\xi|^2\right)
	\]
	which is again finite by Lemma \ref{L:equiv}. Now we may say that that there exists some constant $K_t$ dependent on $t$ such that \eqref{E:fcond} has the following bound which is finite since $0<\alpha <1$
	\begin{align*}
	\eqref{E:fcond} \le K_t\bigg(1+2J_0(t,x)+J_0(t,x)^2\bigg)\int_0^t \ud s \: (t-s)^{\alpha-1}.
	\end{align*}
\end{proof}

\begin{corollary} \label{C:explain}
	Let $u(t,x)$ denote the solution to \eqref{E:spde} and $\alpha \in (0,1)$. Then for $t_0>0$ we may say that \begin{equation}
	u(t_0,x)=\big( G(t_0,\cdot)*\varphi(\cdot) \big)(x)+\dfrac{\sin(\alpha \pi)}{\pi}F_\alpha^{t_0}\big(Y(s,\cdot)\big)(x) \label{E:1fac}.
	\end{equation}
	\begin{proof}
		Follows directly from \eqref{E:fac}.
	\end{proof}
\end{corollary}

We are now able to prove Lemma \ref{L:Main}.


\begin{proof}[Proof of Lemma \ref{L:Main}]
	We first define a set $\mathscr K (\Theta)$ as
	\[
	\mathscr K (\Theta) = \{ \big(G(\delta_0,\cdot)*y(\cdot)\big)(x)+(F_\alpha h)(x),
	\quad \text{where $|y|_{\hat\rho} \le \Theta$  and  $|h|_{L^q(0,1,L^2_{\hat \rho})} \le \Theta $} \}.
	\]
	Then $\mathscr K (\Theta)$ is relatively compact in $L^2_{ \rho}$ for $\Theta >0$. Moreover, by \eqref{E:yineq}, \eqref{E:1fac} and Chebychev's inequality, and if we choose $R$ and $\Theta$ so that $| \hat{\varphi}|_{\hat\rho} < R < \Theta$  then we have
	\begin{align*}
	\mathbb{P}\{ u^{\hat{\varphi}}(\delta_0,\cdot) \not \in \mathscr K(\Theta) \}
	& \le \mathbb{P}\left[ \left(\int_{0}^{1} |Y^{\hat{\varphi}}(s,\cdot)|_{\rho}^q \: \ud s\right)^{1/q} > \dfrac{\pi \Theta}{\sin(\alpha \pi )} \right]
	\\ & \le \dfrac{\sin^q(\alpha \pi)}{\pi^q \Theta^q}\E \int_{0}^{1} |Y^{\hat{\varphi}}(s,\cdot)|_{\hat \rho}^q \: \ud s
	\\ & \le \dfrac{\sin^q(\alpha \pi)}{\pi^q \Theta^q}\bigg( 1+ |{\hat{\varphi}}|_{\hat\rho }^q \bigg)
	\end{align*}
	which implies that
	\begin{align*}
	\{ u^{\hat{\varphi}}(\delta_0,\cdot) \in \mathscr K(\Theta) \}
	& \ge 1- \dfrac{\sin^q(\alpha \pi)}{\pi^q \Theta^q}\bigg( 1+ |{\hat{\varphi}}|_{\hat\rho }^q \bigg)
	\ge 1- \dfrac{\sin^q(\alpha \pi)}{\pi^q \Theta^q}\bigg( 1+ R^q \bigg).
	\end{align*}
	Therefore, if we choose $\Theta>0$ big enough such that $ \dfrac{\sin^q(\alpha \pi)}{\pi^q \Theta^q}\bigg( 1+ R^q \bigg) < \epsilon$ the we get that
	\begin{align*}
	\mathbb{P}\{ u^{\hat{\varphi}}(\delta_0,\cdot) \in \mathscr K(\Theta) \} & \ge 1- \epsilon.
	\end{align*}
	We can now deduce from this and the semigroup property of the solution to the heat equation that
	\begin{align*}
	\mathbb{P}\{ u^{\hat{\varphi}}(t,\cdot) \in \mathscr K(\Theta) \} & \ge (1-
	\epsilon)\mathbb{P}\{ |u^{\hat{\varphi}}(t-\delta_0,\cdot)|_{\hat\rho} < R
	\}.
	\end{align*}
	Since we assumed that $u^{\hat \varphi}$ is bounded in probability in $L^2_{\hat \rho}$ then for and $\epsilon >0$ we can fix $R>0$ such that $\mathbb{P}\{ |u^{\hat{\varphi}}(t,\cdot)|_{\hat\rho} \ge R \} < \epsilon$ and then taking $\Theta > R$ such that $ \dfrac{\sin^q(\alpha \pi)}{\pi^q \Theta^q}\bigg( 1+ R^q \bigg) < \epsilon$ concludes the claim.
\end{proof}

\begin{small}% {{{ References
	\addcontentsline{toc}{section}{Bibliography}
	\begin{thebibliography}{alpha}
		\bibitem[AM03]{AM03}% {{{
		Assing, Sigurd and Ralf Manthey.
		\newblock Invariant measures for stochastic heat equations with unbounded coefficients.
		\newblock {\it Stochastic Process. Appl.} 103 (2003), no. 2, 237--256.
		% }}}
		\bibitem[CM94]{CarMol94}% {{{
		Carmona, Ren\'e~A.  and Stanislav~A. Molchanov.
		\newblock Parabolic Anderson problem and intermittency.
		\newblock {\em Mem. Amer. Math. Soc.}, 108 (1994), no. 518, viii+125 pp.
    % }}}
		\bibitem[Cer01]{Cer01}% {{{
		Cerrai, Sandra.
		\newblock {\it Second order PDE's in finite and infinite dimension.}
		\newblock A probabilistic approach. Lecture Notes in Mathematics, 1762. Springer-Verlag, Berlin, 2001. x+330 pp.
    % }}}
		\bibitem[CD15]{ChenDalang13Heat}% {{{
		Chen, Le and Robert C. Dalang.
		\newblock Moments and growth indices for nonlinear stochastic heat equation with rough initial conditions.
		\newblock {\em Ann.\ Probab.}  43 (2015), no. 6, 3006--3051.
    % }}}
		\bibitem[CHN19]{CHN16Density}% {{{
		Chen, Le, Yaozhong Hu and David Nualart.
		\newblock Regularity and strict positivity of densities for the nonlinear stochastic heat equation.
		\newblock {\em Mem. Amer. Math. Soc.}, 2019, to appear.
    % }}}
		\bibitem[CH19]{CH16Comparison}% {{{
		Chen, Le and Jingyu Huang.
		\newblock Comparison principle for stochastic heat equation on $\R^d$.
		\newblock {\em Ann.\ Probab.}  47 (2019), no. 2, 989--1035.
    % }}}
		\bibitem[CK19]{CK15SHE}% {{{
		Chen, Le and Kunwoo Kim.
		\newblock Nonlinear stochastic heat equation driven by spatially colored noise: moments and intermittency.
		\newblock {\em Acta Math. Sci. Ser. B}, 39B (2019), No. 3, 645--668.
    % }}}
		\bibitem[DK+09]{DalangEtc09Minicourse} % {{{
		Dalang, Robert, Davar Khoshnevisan, Carl Mueller, David Nualart and Yimin Xiao.
		\newblock {\it A minicourse on stochastic partial differential equations}.
		\newblock Held at the University of Utah, Salt Lake City, UT, May 8–19, 2006.  Edited by  Khoshnevisan and Firas Rassoul-Agha. Lecture Notes in Mathematics, 1962.
		\newblock Springer-Verlag, Berlin, 2009. xii+216 pp.
    % }}}
		\bibitem[DQ11]{DalQuer}% {{{
		Dalang, Robert C. and Llu\'is Quer-Sardanyons.
		\newblock Stochastic integrals for spde's: a comparison.
		\newblock {\it Expo. Math.} 29 (2011), no. 1, 67--109.
    % }}}
		\bibitem[DZ92a]{DZ92}% {{{
		Da Prato, Giuseppe and Jerzy Zabczyk.
		\newblock {\it Stochastic equations in infinite dimensions.}
		Encyclopedia of Mathematics and its Applications, 44. Cambridge University Press, Cambridge, 1992. xviii+454 pp.
    % }}}
		\bibitem[DZ92b]{DZp92}% {{{
		Da Prato, Giuseppe and Jerzy Zabczyk.
		\newblock {\it Invariant measures for semilinear stochastic equations.}
		Stochastic Anal. Appl. 10, 387-408.
    % }}}
		\bibitem[DZ96]{DZ96}% {{{
		Da Prato, Giuseppe and Jerzy Zabczyk.
		\newblock {\it Ergodicity for infinite-dimensional systems. }
		\newblock London Mathematical Society Lecture Note Series, 229. Cambridge University Press, Cambridge, 1996.
    % }}}
		\bibitem[Mas99]{MS99}% {{{
		Maslowski, Seidler.
		\newblock {\it On sequentially weakly Fellwe solutions to SPDEs.}
		\newblock Rend. Lincei Mat. Appl. 10, 69-78
    % }}}
		\bibitem[MSY16]{MSY16}% {{{
		Misiats, Oleksandr, Oleksandr Stanzhytskyi, and Nung-Kwan Yip.
		\newblock Existence and uniqueness of invariant measures for stochastic reaction-diffusion equations in unbounded domains.
		\newblock {\it J. Theoret. Probab.} 29 (2016), no. 3, 996--1026.
    % }}}
		\bibitem[MSY15]{MSY15}% {{{
		Misiats, Oleksandr, Oleksandr Stanzhytskyi, and Nung-Kwan Yip.
		\newblock Ito's Formula in infinite dimension and its application to the existence of invariant measures.
    % }}}
		\bibitem{NIST2010}% {{{
		F.~W.~J. Olver, D.~W. Lozier, R.~F. Boisvert, and C.~W. Clark, editors.
		\newblock {\em N{IST} handbook of mathematical functions}.
		\newblock U.S. Department of Commerce National Institute of Standards  and Technology, Washington, DC. Cambridge Univ. Press, Cambridge, 2010.
    % }}}
		\bibitem[TZ98]{TessitoreZabczyk98M}% {{{
		Tessitore, Gianmario and Jerzy Zabczyk.
		\newblock Invariant measures for stochastic heat equations.
		\newblock {\it Probab. Math. Statist.} 18 (1998), no. 2, Acta Univ. Wratislav. No. 2111, 271--287.
    % }}}
		\bibitem[Wal86]{Walsh}% {{{
		Walsh, John B.
		\newblock {\it An Introduction to Stochastic Partial Differential Equations}.
		\newblock In: \`Ecole d'\`et\'e de probabilit\'es de
		Saint-Flour, XIV---1984, 265--439.
		Lecture Notes in Math.\ 1180, Springer, Berlin, 1986.
    % }}}
	\end{thebibliography}
\end{small}
% }}}
\end{document}
